% !TEX TS-program = pdflatex
% !BIB program = bibtex
\documentclass[11pt,a4paper]{article}
% Shared packages, typography and course navigation.
\usepackage[T1]{fontenc}
\usepackage[utf8]{inputenc}
\usepackage{lmodern,microtype}
\usepackage[margin=25mm]{geometry}
\usepackage{amsmath,amssymb,graphicx,booktabs,tabularx,array,longtable}
\usepackage[round,authoryear]{natbib}
\usepackage[dvipsnames]{xcolor}
\usepackage{tcolorbox}
\tcbuselibrary{breakable}
\usepackage{tikz}
\usetikzlibrary{arrows.meta,positioning}
\usepackage{caption,placeins,fancyhdr}
\usepackage[colorlinks=true,linkcolor=MidnightBlue,citecolor=MidnightBlue,urlcolor=MidnightBlue]{hyperref}
\urlstyle{same}
\setlength{\emergencystretch}{3em}
\setlength{\parskip}{5pt}
\setlength{\parindent}{0pt}
\renewcommand{\arraystretch}{1.22}
\captionsetup{font=small,labelfont=bf}
\pagestyle{fancy}\fancyhf{}\fancyhead[L]{\small From Brains to Models}\fancyfoot[C]{\thepage}
\setlength{\headheight}{14pt}
\newcommand{\learningcheck}[2]{%
  \par\smallskip\noindent\begin{minipage}{\linewidth}
  \textbf{Takeaway.} #1\par\smallskip
  \textbf{Check your understanding.} #2\end{minipage}\par\smallskip}
\newcommand{\coursenotebook}[2]{\href{https://github.com/arvidl/brains-to-models/blob/main/code/notebooks/#1}{#2}}
\newtcolorbox{perspectives}[1]{%
  title={Two reading paths: #1},
  colback=blue!2,colframe=MidnightBlue!45,
  colbacktitle=blue!6,coltitle=MidnightBlue,
  fonttitle=\small\bfseries,fontupper=\small,
  boxrule=0.4pt,arc=1mm,
  left=2.5mm,right=2.5mm,top=1.5mm,bottom=1.5mm,
  before skip=8pt,after skip=8pt}

% Stable links between independently compiled course documents.
\newcommand{\coursepaper}[1]{\href{https://github.com/arvidl/brains-to-models/blob/main/paper/main.pdf}{#1}}
\newcommand{\coursesyllabus}[1]{\href{https://github.com/arvidl/brains-to-models/blob/main/paper/syllabus.pdf}{#1}}
\newcommand{\coursesolutions}[1]{\href{https://github.com/arvidl/brains-to-models/blob/main/paper/syllabus-solutions.pdf}{#1}}
\newcommand{\courseauthor}{Arvid Lundervold\\\normalsize Department of Biomedicine, University of Bergen\\\normalsize Mohn Medical Imaging and Visualization Centre}
\newcommand{\coursedate}{\small Companion to the Joint PRESIMAL $\times$ NORA HS keynote\\5 November 2026, Sola Strand Hotel, Stavanger}

\newcommand{\currentsolutionheading}{Tentative syllabus solutions}
\fancyhead[R]{\small\currentsolutionheading}
\hypersetup{pdftitle={Tentative solutions to the syllabus exercises},pdfauthor={Arvid Lundervold}}
\title{\textbf{Tentative solutions to the\\syllabus exercises}\\[5mm]\large From Brains to Models: What AI Reveals About Being Human}
\author{\courseauthor}
\date{\coursedate}
\newcounter{solutionslide}
\newcommand{\solutionslide}[2]{%
  \clearpage\refstepcounter{solutionslide}
  \gdef\currentsolutionheading{Slide #1: Tentative solutions}
  \section*{Slide #1. #2}\label{solution-#1}
  \addcontentsline{toc}{section}{Slide #1. #2}}
\newcommand{\biomedicalanswer}{\pdfbookmark[1]{Biomedical and imaging path}{answer-bio-\arabic{solutionslide}}\subsection*{Biomedical and imaging path}}
\newcommand{\computationalanswer}[1]{\clearpage\pdfbookmark[1]{Computational path}{answer-comp-#1}\subsection*{Slide #1: Computational path}}
\newcommand{\answercontinued}[1]{\clearpage\subsection*{Slide #1: Computational path (continued)}}
\newcommand{\answerstep}[1]{\par\medskip\textbf{#1}\par\nobreak}
\newcommand{\answerconnection}[1]{\par\medskip\textbf{Connecting the paths.} #1\par}
\newcommand{\answercheck}[1]{\par\medskip\textbf{Check your reasoning.} #1\par}
\newtcolorbox{exercisefocus}{colback=blue!2,colframe=MidnightBlue!35,boxrule=.4pt,arc=1mm,left=2.5mm,right=2.5mm,top=1.5mm,bottom=1.5mm,fontupper=\small,before skip=7pt,after skip=9pt,breakable}
\newcommand{\task}[1]{\begin{exercisefocus}\textbf{Exercise in focus.} #1\end{exercisefocus}}
\begin{document}
\maketitle
\section*{How to use these tentative solutions}
This separate companion develops answers to the exercises for all 25 slides in the \coursesyllabus{slide-by-slide syllabus}. Attempt a question first, then compare the reasoning, assumptions and conclusion. Slide numbers and titles match the syllabus; each slide's biomedical answer and computational answer begin on separate pages. Longer derivations continue as needed.

\textbf{Two ways of reasoning.} The \emph{Biomedical and imaging path} starts with a phenomenon, explains it in ordinary biological or clinical language, considers alternatives and proposes an informative observation. Open questions admit several defensible answers; the examples here show what makes one answer persuasive. The \emph{Computational path} specifies mathematical objects and units, works through a calculation or proof and interprets precisely what follows. Its conclusions are conditional on its assumptions. A brief connection brings the two results together.

\textbf{What ``tentative'' means.} These are worked teaching answers for discussion and revision, rather than a unique marking key. Algebraic results can be checked exactly; a proposed clinical interpretation or study design needs further evidence. Patient situations and numerical examples are hypothetical unless explicitly identified as saved course results. Suggested notebook experiments describe what to run and record; this document does not present them as newly executed experiments.

\textbf{Read symbols locally.} Each computational answer reintroduces its essential notation. A symbol can change meaning between examples: $D$ may be a dataset, diffusivity or disease event, and $p$ may be an event probability or a perfusion factor. The syllabus's glossary, before its references, supplies abbreviations, biological concepts and mathematical notation for both paths. The \coursepaper{perspective paper} provides the broader discussion and technical appendices; notebook links lead to the executable examples.

\clearpage
\begingroup\small
\makeatletter\let\l@section\l@subsection\makeatother
\tableofcontents\endgroup
\solutionslide{01}{From brains to models}
\biomedicalanswer
\task{Choose a lesion-follow-up task. Explain what the contour means, what might mislead it and what a patient would need from its use.}
\answerstep{1. Give the contour a specific meaning.}
Consider following a region of abnormal signal on successive brain MRI examinations. A contour marks the region included by an annotation policy on those images. For example, a policy may include a specified FLAIR abnormality. It does not automatically mark every tumour cell, identify a disease mechanism or establish that the person needs a particular treatment. The biological meaning depends on which image contrast was outlined and on the clinical context.

A useful ordinary-language explanation is: ``This outline helps us compare the visible region across scans. It is a measurement of an image appearance. We still need to interpret what a change means in your situation.'' That explanation separates a reproducible measurement from a diagnosis or decision.
\answerstep{2. Explain a plausible misleading agreement.}
A model and a radiologist might draw similar contours on a familiar scan, while the model relies heavily on brightness. A change in acquisition, tissue fluid or the boundary convention could then alter the outline without a corresponding change in viable tumour. Conversely, two methods might agree because both use the same imperfect reference policy. Agreement is useful evidence about the specified output, but does not identify the reasoning that produced it.
\answerstep{3. State what would make the output useful.}
The patient would need the measurement related to symptoms, prior scans, uncertainty and the purpose of follow-up. The team should ask whether the apparent change exceeds plausible measurement variation and whether using it improves a meaningful part of care. Reproducibility, interpretation and benefit are separate questions.
\answercheck{A satisfactory answer names the outlined contrast and purpose, offers an alternative explanation for a changed contour and identifies what additional information would matter. ``The model agrees with an expert, therefore it understands the disease'' skips those steps.}

\computationalanswer{01}
\task{Construct two real-valued functions agreeing on $D=\{0,1\}$ but differing at $x=2$. Explain what the additional test establishes.}
\answerstep{1. Define the functions and evaluate them.}
Let the input and output sets both be $\mathbb R$, the real numbers. Choose
\[
f(x)=0,\qquad g(x)=x(x-1).
\]
These are functions defined for every real input, not just for the two observed cases. Direct substitution gives
\[
\begin{array}{c|ccc}
x&0&1&2\\\hline
f(x)&0&0&0\\
g(x)&0&0&2
\end{array}
\]
Thus the agreement fraction on $D$ is
\[
A_D(f,g)=\tfrac12\bigl(\mathbf1[0=0]+\mathbf1[0=0]\bigr)=1.
\]
The indicator $\mathbf1[Q]$ equals one when statement $Q$ is true. Perfect agreement on the finite set does not imply the universal statement $f(x)=g(x)$ for every real $x$.
\answerstep{2. Add one observation and state its consequence.}
If the newly observed target at $x=2$ is zero, $g$ is inconsistent with the three exact observations and $f$ remains consistent. This rejects this particular $g$ as an exact predictor of all three cases. It does not prove that $f$ is the unique possible rule. For any nonzero constant $a$,
\[
h_a(x)=a\,x(x-1)(x-2)
\]
also returns zero at all three tested points but differs elsewhere. The factors make the construction easy to read: each factor forces a zero at one observed input.
\answerstep{3. Generalize and interpret.}
For any finite collection of distinct inputs $x_1,\ldots,x_n$, the function
\[
g_a(x)=f(x)+a\prod_{i=1}^n(x-x_i)
\]
agrees with $f$ on those inputs. The product vanishes whenever $x=x_i$. Without further restrictions on the admissible function family, finite observations leave behaviour elsewhere undetermined.

Even identifying a complete input--output function would not uniquely identify an implementation: a formula, a table on a restricted domain and different algorithms can produce the same answers. Mechanism claims therefore require observations of internal structure, dynamics or responses to interventions in addition to output agreement. \coursenotebook{00_intro_timeline.ipynb}{Notebook 00} introduces this separation through Marr's levels \citep{marr2010}.
\answerconnection{The biomedical answer specifies why a contour matters. The counterexample establishes exactly why matching contours on tested cases leaves both untested behaviour and internal mechanisms open.}

\solutionslide{02}{A learning programme}
\biomedicalanswer
\task{Make a verbal prediction before a notebook run, explain a possible surprise and name an observation that could resolve it.}
\answerstep{1. Choose a prediction that can be checked.}
For \coursenotebook{15_ppv_and_decisions.ipynb}{Notebook 15}, a learner might predict: ``If a test detects nine out of ten affected people and correctly clears nine out of ten unaffected people, most positive results will indicate disease.'' This is a plausible starting intuition, not the correct conclusion for every population.

Before running anything, specify the missing context: suppose only one in a hundred people has the condition. A prediction about positive results cannot be assessed without this population information. Write the prediction down so it is not retrospectively changed to match the output.
\answerstep{2. Explain the surprise using people rather than formulas.}
In the syllabus's hypothetical population of 10,000, 100 people have disease and 9,900 do not. Expect 90 true positives among the first group and 990 false positives among the second. The large unaffected group supplies many more positive results, despite its lower chance of testing positive. Consequently, only about eight of every hundred positive results correspond to disease in this example.

The surprise is explained by the relative sizes of the two groups. It does not require the test to change between them. These are expected counts from stipulated inputs; a real sample will fluctuate.
\answerstep{3. Name the observation needed outside the notebook.}
To investigate whether the explanation applies to a proposed setting, estimate condition frequency in the relevant population and compare test results with an appropriate independent reference for condition status. Record who was included and how the reference was obtained. A conveniently selected group may not represent people who would actually be screened.
\answercheck{Running the notebook checks what the stipulated inputs imply. It cannot by itself establish those inputs for a clinical population. A useful explanation names which assumption caused the surprise and which observation could challenge that assumption.}

\computationalanswer{02}
\task{Give two populations compatible with the same finite training sample but yielding different risks for a fixed predictor. Explain the limit of sample-loss minimization.}
\answerstep{1. Specify the predictor, targets and loss.}
Let $\mathcal X=\{a,b\}$ and $\mathcal Y=\{0,1\}$. Suppose the true target is zero at either input. Fix a predictor with $f(a)=0$ and $f(b)=1$. Use zero--one loss,
\[
L(\hat y,y)=\mathbf1[\hat y\ne y].
\]
This loss is dimensionless. The predictor makes no error at $a$ and an error at $b$.
\answerstep{2. Define two possible joint distributions.}
Under $P_1$, let $P_1(a,0)=0.9$ and $P_1(b,0)=0.1$. Under $P_2$, reverse these probabilities. All other input--target pairs have probability zero. The expected, or population, risks are
\[
\begin{aligned}
R_{P_1}(f)&=0.9\times0+0.1\times1=0.1,\\
R_{P_2}(f)&=0.1\times0+0.9\times1=0.9.
\end{aligned}
\]
The two populations therefore give very different error probabilities to the same function. The changed object is $P$, not $f$ or $L$.
\answerstep{3. Check compatibility with an observed sample.}
Consider $n>0$ independent training observations, all equal to $(a,0)$. This sample has probability $0.9^n$ under $P_1$ and $0.1^n$ under $P_2$. Both are positive, so the same observed finite sample is possible under either population. The sample is more likely under one population; the exercise does not require equal likelihood.

For that sample,
\[
\widehat R_n(f)=\frac1n\sum_{i=1}^n L(f(a),0)=0.
\]
Zero is the smallest possible sample loss, so $f$ is a sample minimizer if included in the candidate family. Its population risk nevertheless depends on the unobserved mixture of $a$ and $b$.
\answerstep{4. Interpret the distinction.}
Sample minimization establishes an optimization property on observed cases. Generalization requires assumptions connecting the sample and target population, together with appropriate evaluation. A demonstration shows that a computation is possible; verification checks its implementation against the specified mathematics; validation asks whether the model is adequate against independent observations for a stated use. The three claims have different evidence requirements.
\answerconnection{Population counts make the prevalence effect intuitive. The construction formalizes why an observed training average cannot determine every relevant population average.}

\solutionslide{03}{Biological brains across scales}
\biomedicalanswer
\task{Describe a visible contrast, propose a cellular explanation and an alternative, and identify an additional measurement that distinguishes them.}
\answerstep{1. Keep description separate from explanation.}
A defensible visual description is that the drawings show different arrangements of cell bodies and branching processes, while the photographs show larger-scale differences in brain form. Some illustrated regions appear more densely filled with processes than others. These statements concern what is visible. They do not establish which brain is more intelligent or how much tissue a drawing represents. Scale, section orientation and preparation matter.
\answerstep{2. Offer two explanations for one contrast.}
One possible explanation for a denser-looking microscopic panel is a difference in the number or branching pattern of labelled cells. An alternative is a difference in section thickness, staining, sampling location or the direction in which processes cross the section. The same apparent density can arise through either biological organization or how it was measured.

An analogous imaging example is that one MRI appearance can be compatible with several microscopic states. The signal aggregates contributions within a voxel; it does not count a named cell population directly. Calling the appearance a measurement of ``more cells'' already adds a hypothesis.
\answerstep{3. Choose a discriminating observation.}
To compare microscopic branching, use matched tissue preparation, calibrated spatial scales and a defined sampling scheme. Three-dimensional reconstructions or an independent cell-counting measurement could help distinguish increased branching from more labelled cells. The additional observation should address the competing explanation, rather than repeat the same ambiguous image.

For an imaging hypothesis, specify how the proposed additional tissue or imaging measurement relates to the hidden state. More observations help only if they contain information not already lost by the first measurement.
\answercheck{A strong answer states what the image supports before naming a mechanism. It treats scale bars and measurement preparation as part of the evidence and explains why the proposed second measurement would distinguish the alternatives.}

\computationalanswer{03}
\task{Recover a two-component state from its sum and an added first-component measurement. Identify architecture, encoding and update-rule differences in Notebook 03.}
\answerstep{1. Show why the first observation is insufficient.}
Let $s=(s_1,s_2)^\top$ be a column vector of two state components. The noiseless scalar observation is
\[
x=Hs,\qquad H=\begin{bmatrix}1&1\end{bmatrix}.
\]
For example, $s=(1,2)^\top$ and $s=(2,1)^\top$ both give $x=3$. The vector $v=(1,-1)^\top$ satisfies $Hv=0$, so adding any multiple of $v$ to a state preserves its measured sum. This invisible direction is the null space of the observation map.
\answerstep{2. Add a genuinely independent observation.}
Measure $r=s_1$ in addition to $x=s_1+s_2$. The combined system is
\[
\begin{bmatrix}x\\r\end{bmatrix}
=\underbrace{\begin{bmatrix}1&1\\1&0\end{bmatrix}}_{\widetilde H}
\begin{bmatrix}s_1\\s_2\end{bmatrix}.
\]
The determinant of $\widetilde H$ is $-1$, which is nonzero; the two rows supply independent linear constraints. Recovery is particularly transparent without matrix inversion:
\[
s_1=r,\qquad s_2=x-r.
\]
Thus $x=3,r=1$ recover $(1,2)$, whereas $x=3,r=2$ recover $(2,1)$. Measuring the sum again would not remove the ambiguity.
\answerstep{3. State the assumptions behind recovery.}
Exact recovery here assumes the stated measurement maps, the same underlying state for both measurements and no measurement error. With errors $\epsilon_x,\epsilon_r$, the reconstructed components become $s_1+\epsilon_r$ and $s_2+\epsilon_x-\epsilon_r$. Uniqueness of the noiseless solution therefore does not imply zero uncertainty in practice. Units must also agree before the subtraction is meaningful.
\answerstep{4. Identify the notebook's distinct design choices.}
\coursenotebook{03_bio_vs_artificial.ipynb}{Notebook 03} compares a spiking perceptron with a multilayer perceptron on shared constructed tasks. Architecture differs: a shallow spiking decision system versus a network with a hidden layer. Encoding differs: Poisson input spikes versus numerical input features. Learning differs: a local update using pre-synaptic activity and an output error versus error propagation through the hidden layer by gradient-based training.

These changes occur together. An accuracy difference cannot isolate the effect of locality, spiking or depth. A more discriminating comparison holds architecture and encoding fixed while changing the update rule, or explicitly studies their interactions.
\answerconnection{The biological answer asks which observation would separate mechanisms. The linear example explains what ``separate'' requires: a second measurement must constrain a direction invisible to the first.}

\solutionslide{04}{Architectures encode choices}
\biomedicalanswer
\task{Explain how a model can perform well on large lesions but miss small ones, and identify an observation hidden by an average score.}
\answerstep{1. Explain a plausible size-dependent failure.}
A large lesion occupies many image elements and often remains visible after spatial downsampling. A small lesion may occupy only a few elements, share a voxel with surrounding tissue or disappear from a coarser internal representation. A boundary error of a few voxels can also represent a large fraction of a small lesion. These are plausible mechanisms to investigate, rather than a diagnosis of every segmentation architecture.

The task matters: outlining an already identified large region differs from detecting every small separate focus. A system may succeed at the former while failing at the latter. U-Net and attention-based architectures organize information differently, but their names do not establish detection performance \citep{ronneberger2015,vaswani2017}.
\answerstep{2. Construct an average that hides the problem.}
Imagine two reference lesions: one contains 1,000 voxels and the other 10. Suppose the model outlines the large lesion perfectly and misses the small lesion completely. The pooled Dice overlap is $2,000/2,010\approx0.995$. Yet it detected only one of the two lesions. If Dice were averaged separately over these two lesions, the result would instead be $(1+0)/2=0.5$ under this simple matching convention.

All these numbers describe the same invented masks. They differ because they weight the errors differently.
\answerstep{3. Name the missing observation.}
Report lesion detection by size and location, the number of missed separate foci, and the clinical relevance of those omissions under the intended task. Inspect individual cases as well as pooled overlap. A high average contour score can conceal an entirely absent lesion that matters more than a small boundary discrepancy on a large one.
\answercheck{Explain how image scale and the averaging rule contribute separately. ``The average is high'' does not answer whether the clinically important small lesions were found.}

\computationalanswer{04}
\task{Check attention-array dimensions. Keep attention weights $A$ fixed, replace values $V$ by $2V$, and calculate the resulting output.}
\answerstep{1. Check that the products are defined.}
For $n$ tokens, let $Q,K\in\mathbb R^{n\times d_k}$ and $V\in\mathbb R^{n\times d_v}$. The transpose $K^\top$ has $d_k$ rows and $n$ columns, so
\[
QK^\top\in\mathbb R^{n\times n},\quad
A=\operatorname{softmax}(QK^\top/\sqrt{d_k})\in\mathbb R^{n\times n},\quad
Z=AV\in\mathbb R^{n\times d_v}.
\]
For scaled score $s_{ij}$, row-wise softmax defines
\[
A_{ij}=\frac{e^{s_{ij}}}{\sum_{r=1}^n e^{s_{ir}}}.
\]
Every term is positive and summing over column $j$ returns one. Row $i$ supplies weights for combining the value vectors; it is not itself that combined vector.
\answerstep{2. Work through a two-token example.}
Choose $n=2$, $d_v=1$ and
\[
A=\begin{bmatrix}3/4&1/4\\1/2&1/2\end{bmatrix},\qquad
V=\begin{bmatrix}1\\3\end{bmatrix}.
\]
Then
\[
Z=AV=\begin{bmatrix}(3/4)1+(1/4)3\\(1/2)1+(1/2)3\end{bmatrix}
=\begin{bmatrix}1.5\\2\end{bmatrix}.
\]
Changing the values to $2V=(2,6)^\top$ gives
\[
Z'=A(2V)=2(AV)=\begin{bmatrix}3\\4\end{bmatrix}.
\]
The equality follows from linearity of matrix multiplication when $A$ is fixed. Exactly the same attention weights can therefore produce different outputs.
\answerstep{3. Interpret the scope of this result.}
The result concerns one attention operation, not every component of a Transformer. Doubling the original input can change $Q,K$ and hence $A$; projections, biases and normalization can also change the full model's response. Holding $A$ fixed is an explicit experimental condition.

In \coursenotebook{01_tiny_transformer.ipynb}{Notebook 01}, trace score formation, normalization and value aggregation separately. In \coursenotebook{09_segmentation_reflection.ipynb}{09}, distinguish the output scores from the binary mask formed by a threshold. Neither an attractive attention map nor a score bounded by zero and one alone establishes a clinical explanation or calibrated probability.
\answerconnection{The biomedical failure concerns whether relevant spatial information reaches the output. The calculation identifies what an attention map can and cannot specify without the associated representations.}

\solutionslide{05}{Four questions about intelligence}
\biomedicalanswer
\task{Place four questions about one diagnostic decision in the human/rational and thought/action grid. Explain the evidence each question requires.}
\answerstep{1. Keep one decision fixed across the grid.}
Consider an uncertain image finding and the question of whether to request further assessment. The same setting supports four different investigations:
\begin{itemize}
\item \textbf{Thinking humanly:} Does the system use a process resembling how a clinician reasons? Relevant evidence could include the sequence of information sought, responses to conflicting findings and error patterns. A fluent explanation alone does not reveal the actual process.
\item \textbf{Acting humanly:} Does its recommendation resemble those made by clinicians under the same information? Compare recommendations on matched cases. Agreement can reproduce shared habits or errors.
\item \textbf{Thinking rationally:} Does the conclusion follow from the stated premises and rules? Check logical validity and then separately ask whether the clinical premises are justified.
\item \textbf{Acting rationally:} Does the selected action do best under declared probabilities and consequences? Specify which outcomes matter, to whom, and how uncertainty and burden are represented.
\end{itemize}
\answerstep{2. Show why one success can accompany another failure.}
A model could reproduce a clinician's usual recommendation while relying on an irrelevant image marker. It would agree behaviourally without establishing a similar reasoning mechanism. Alternatively, an impeccably valid argument starting from ``every enhancing region is viable tumour'' could be unhelpful because the premise is too strong. Agreement and formal validity do not repair an unsupported clinical assumption.

A recommendation preferred under one person's priorities may also change when the burdens or consequences are valued differently. A formal action criterion makes those supplied values visible; it does not settle them on the patient's behalf.
\answercheck{The four questions should concern the same task but require different evidence. Treating a human-like explanation as simultaneous proof of mechanism, truth and usefulness collapses the grid.}

\computationalanswer{05}
\task{Construct two policies ranked differently by agreement with a human and by expected utility. Supply the probabilities, utilities and calculations.}
\answerstep{1. Define a deliberately simple decision problem.}
Let outcome $Y\in\{0,1\}$ denote absence or presence of an event. Given the available evidence $x$, stipulate $P(Y=1\mid x)=0.1$. The two actions are $a_{\rm act}$ and $a_{\rm wait}$. Define utility as negative loss:
\[
\begin{array}{c|cc}
U(a,y)&y=0&y=1\\\hline
a_{\rm act}&-1&0\\
a_{\rm wait}&0&-4
\end{array}
\]
The numbers are invented, expressed in a common utility scale. Correct actions cost zero, acting unnecessarily costs one and failing to act when the event occurs costs four.
\answerstep{2. Compute expected utilities.}
Expected utility weights each possible consequence by its conditional probability:
\[
\begin{aligned}
\mathbb E[U(a_{\rm act},Y)\mid x]&=0.9(-1)+0.1(0)=-0.9,\\
\mathbb E[U(a_{\rm wait},Y)\mid x]&=0.9(0)+0.1(-4)=-0.4.
\end{aligned}
\]
Because $-0.4>-0.9$, waiting maximizes expected utility under these stipulated inputs.
\answerstep{3. Define agreement and obtain the reversed ranking.}
Suppose the reference human chooses $a_{\rm act}$ on every case in a set of cases with these same probabilities and costs. Policy $\pi_1$ always acts; policy $\pi_2$ always waits. Their agreement fractions with this reference are
\[
A(\pi_1)=1,\qquad A(\pi_2)=0.
\]
Agreement ranks $\pi_1$ above $\pi_2$, while expected utility ranks $\pi_2$ above $\pi_1$. This counterexample disproves the general implication ``greater human agreement means greater expected utility.'' It does not estimate how often clinicians make either choice.
\answerstep{4. Interpret the formal criteria.}
If the event probability were instead 0.3 with the same costs, acting would have utility $-0.7$ and waiting $-1.2$; the preferred action would reverse. Utility conclusions depend on the supplied probability and value model. Likewise, $\Gamma\vdash\phi$ means derivability of conclusion $\phi$ from premises $\Gamma$, not empirical truth of $\Gamma$. \coursenotebook{04b_four_views_of_AI.ipynb}{Notebook 04b} separates these targets.
\answerconnection{Clinical purposes determine which success question matters. The counterexample demonstrates mathematically why success under one criterion cannot substitute for the others.}

\solutionslide{06}{Embodiment, judgment and collaboration}
\biomedicalanswer
\task{Explain how fatigue could change a choice while the image stays identical. Name a bodily observation and an alternative explanation.}
\answerstep{1. Describe a plausible change in the person.}
The same scan can be encountered by a person in different bodily states. A fatigued reader might find sustained search more demanding, seek an additional opinion or stop examining an ambiguous boundary sooner. Another reader might compensate by becoming more cautious. Fatigue need not have one universal behavioural direction. The point is that unchanged external pixels do not imply an unchanged perceptual and decision situation.

Interoception concerns signals about internal bodily condition; homeostatic regulation concerns maintaining physiological variables within viable ranges. These concepts motivate questions about how bodily needs influence cognition \citep{seth2013,sethbayne2022}. They do not make every performance change evidence about consciousness.
\answerstep{2. Propose observations that can separate explanations.}
Record sleep history and a declared fatigue rating alongside the task. A bodily measure related to arousal might provide an additional observation, but no single measure should be treated as a complete definition of fatigue. Record the choice, search duration, confidence and relevant accuracy outcome separately.

An alternative explanation is that the second encounter occurs under a different workload, incentive or instruction. Familiarity with the case can also change the answer. A within-person comparison should therefore consider order, prior exposure and matched task conditions; simply showing the same image twice does not isolate a bodily cause.
\answerstep{3. State the interpretation at the right level.}
If measured fatigue and choices change together, that is an association to explain. A controlled comparison can strengthen a specific causal claim, but the result remains about the measured state and behaviour. It does not establish that a software controller has bodily experience, nor that the reader's complete internal state has been measured.
\answercheck{A good answer names an observable bodily correlate, a specific choice and a credible competing explanation. It distinguishes a hypothesis about bodily influence from a claim to have measured subjective experience.}

\computationalanswer{06}
\task{Establish invariance for an unlimited controller, then show why bounded actions can fail when disturbances are larger.}
\answerstep{1. Define the state, action and disturbance.}
The scalar model is $s_{t+1}=s_t+a_t+w_t$. All quantities here are dimensionless. The viable set is $\mathcal V=[-1,1]$. Full observation means the controller knows $s_t$ before choosing $a_t$. Invariance requires the next state to remain in $\mathcal V$ for every allowed current state and disturbance under the declared policy.
\answerstep{2. Prove the unlimited-action result.}
Use policy $a_t=-s_t$. Substitution gives
\[
s_{t+1}=s_t-s_t+w_t=w_t.
\]
For every $|w_t|\leq0.2$, it follows that $|s_{t+1}|\leq0.2\leq1$. Starting from any $s_0\in\mathcal V$, apply this argument at each step: the next state is in $\mathcal V$, so the same guarantee applies again. This is an induction over updates.

The proof depends on exact observation, the stated transition equation, timely control and permission to choose $-s_t$. It does not cover an unmodelled delay or action constraint.
\answerstep{3. Construct a failure that no allowed action repairs.}
Now impose $|a_t|\leq0.5$ and allow $|w_t|\leq0.6$. Choose the allowed boundary state $s_t=1$ and disturbance $w_t=0.6$. The most negative allowed action is $-0.5$. Consequently, for every admissible action,
\[
s_{t+1}=1+a_t+0.6\geq1-0.5+0.6=1.1>1.
\]
The next state must leave $\mathcal V$. This is stronger than showing that one poorly chosen action fails: even the best allowed correction is insufficient for this state--disturbance pair. At the negative boundary, $s_t=-1,w_t=-0.6$ gives the symmetric failure.
\answerstep{4. Read the quantifiers and their meaning.}
A robust invariance claim has the structure ``for every allowed state, the policy succeeds for every allowed disturbance.'' One permitted pair that defeats all permitted actions is enough to refute that claim. The successful unlimited controller and failed bounded controller have different assumptions, so their conclusions are compatible.

This calculation establishes a functional property of a specified control loop. Relating such a loop to biological regulation is a further modelling question; attributing experience requires another argument. Appendix C of the \coursepaper{paper} develops that distinction.
\answerconnection{Bodily regulation motivates the viable range. The proof clarifies which sensing, action and disturbance assumptions actually guarantee staying within it.}

\solutionslide{07}{What a Turing-like test measures}
\biomedicalanswer
\task{Explain how a scan explanation could sound human while omitting important uncertainty. Propose an observation of comprehension rather than resemblance.}
\answerstep{1. Construct the communication problem.}
Suppose the source finding is an abnormal region whose cause remains uncertain. A polished explanation says, ``The scan shows the cause of your symptoms, and the next step is clear.'' Its familiar phrasing and confident narrative may sound like a human conversation. However, it omits the source's uncertainty and introduces a connection to symptoms that the imaging observation alone has not established.

The resemblance judgment asks who appears to have produced the answer. The comprehension question asks what the recipient understands about the finding. A response can succeed on the first and fail on the second. The syllabus's Turing-like comparison concerns a specified judgment task, not a complete assessment of clinical communication \citep{zhang2026turing}.
\answerstep{2. Define a comprehension observation.}
Ask the recipient: ``What has this scan shown, and what remains uncertain?'' A meaningful answer might distinguish a visible abnormality from its possible causes and recognize that other information is needed. An answer such as ``The scan has proved the cause'' would reveal the particular misunderstanding in this hypothetical example.

Use a predeclared scoring guide based on the source finding. Assess the recipient's own explanation, rather than asking only whether the wording was reassuring or sounded professional. Confidence in an answer and comprehension of its uncertainty are different outcomes.
\answerstep{3. Design a fair comparison.}
Compare explanations for matched source findings, keep essential clinical content explicit and have assessors score comprehension without knowing which system produced the wording. Include recipients with relevant language and health-literacy backgrounds. Record resemblance, comprehension and preference separately so that one outcome cannot silently substitute for another.
\answercheck{An answer should identify the omitted uncertainty, the mistaken belief it could produce and an observable response that would reveal that belief. Human-like style alone supplies none of those measurements.}

\computationalanswer{07}
\task{Derive the variance of an average of independent Bernoulli judgments and of duplicated copies of one judgment. Identify the independent sampling units.}
\answerstep{1. Derive the variance of one correctness indicator.}
Let $Z\in\{0,1\}$ indicate a correct judgment, with $P(Z=1)=a$. Then $Z^2=Z$, so
\[
\mathbb E Z=a,\qquad
\operatorname{Var}(Z)=\mathbb E Z^2-(\mathbb E Z)^2=a-a^2=a(1-a).
\]
Here $a$ is an accuracy probability and is unrelated to an action symbol on other slides.
\answerstep{2. Average independent judgments.}
For $N>0$ independent, identically distributed indicators $Z_1,\ldots,Z_N$, define $\hat a=N^{-1}\sum_k Z_k$. The general variance identity is
\[
\operatorname{Var}(\hat a)=\frac1{N^2}
\left(\sum_k\operatorname{Var}(Z_k)+2\sum_{k<l}\operatorname{Cov}(Z_k,Z_l)\right).
\]
Independence sets the covariances to zero. The $N$ equal variance terms then give $\operatorname{Var}(\hat a)=a(1-a)/N$. Its square root is the standard error, the spread of the estimated accuracy over repeated samples under these assumptions.
\answerstep{3. Replace independent observations by copies.}
If every row is the same random variable $Z$, then
\[
\hat a=\frac1N\underbrace{(Z+\cdots+Z)}_{N\text{ copies}}=Z,
\qquad \operatorname{Var}(\hat a)=a(1-a).
\]
Equivalently, each covariance equals $\operatorname{Var}(Z)$, leaving $N^2$ equal contributions before division by $N^2$. No new information was created by duplicating rows.

For the illustrative value $a=0.6$, one judgment has variance $0.24$ and standard deviation about $0.490$. One hundred independent judgments have mean variance $0.0024$ and SE about $0.049$. One hundred copies of one judgment retain variance $0.24$.
\answerstep{4. Map the mathematics to a proposed study.}
Judgments may share a judge, source item or generated response. Those shared components can induce dependence without making every judgment an exact copy. Identify the sampled judges and items, which responses were reused, and whether conclusions concern these items, new items, these judges or a wider population. Analysis may need to represent both judge and item clustering. The raw number of rows alone does not justify the independent-Bernoulli formula.
\answerconnection{A human resemblance judgment is meaningful only in its task and population. The covariance calculation explains why repeated judgments also need a clear sampling structure.}

\solutionslide{08}{Read results in their time and protocol}
\biomedicalanswer
\task{Describe a presentation change that could alter a patient's impression without improving clinical content, and identify an observation that would detect it.}
\answerstep{1. Hold the informational content fixed.}
Consider two versions of the same explanation: identical statements about the finding, alternatives and uncertainty, presented either as plain text or as an attractive narrated display. A confident voice, familiar layout or animation could change perceived authority, warmth or reassurance without adding a clinical fact. The direction and size of this effect are empirical questions, not assumed results.

A different example is replacing cautious visual emphasis with a vivid highlighted contour while keeping the words unchanged. The appearance may imply a more definite boundary than the accompanying explanation supports.
\answerstep{2. Measure impression and understanding separately.}
Ask recipients to rate perceived confidence or trust, and separately ask them to explain what is known and uncertain. Check factual content against the source using a fixed rubric. A changed impression alongside unchanged factual content and comprehension would support the proposed distinction. Improved comprehension would be a separate possible benefit to measure, rather than assume.

If the same participant sees both versions, order and memory may affect the result. Separate assigned groups or a counterbalanced design can help address that issue, depending on the question and population.
\answerstep{3. Preserve the historical comparison.}
A result obtained for one model version, interface and set of instructions belongs to that encounter. A later model with tools, a different prompt or another display is a changed procedure. Report those changes before interpreting differences. The source study discussed in the syllabus should therefore be read with its model versions and comparison protocol in view \citep{zhang2026turing}.
\answercheck{The answer should identify which presentation feature changed, demonstrate that the clinical propositions stayed fixed and use different observations for impression and comprehension. A preference rating alone cannot establish improved clinical content.}

\computationalanswer{08}
\task{Separate item variability from generation variability, explain protocol confounding and construct a scoring change that reverses a model ranking.}
\answerstep{1. Index the two sources of variability.}
Let $S_{jk}$ be the score on item $j$ in repeated generation $k$. A useful decomposition is $S_{jk}=\mu_j+\epsilon_{jk}$, where $\mu_j$ is that item's expected score and the residual has conditional mean zero. The law of total variance gives
\[
\operatorname{Var}(S)=\operatorname{Var}(\mathbb E[S\mid j])+
\mathbb E[\operatorname{Var}(S\mid j)].
\]
The first term concerns differences between items; the second concerns repeated responses within an item. Many generations of one easy item do not measure variation across a population of difficult and easy items.
\answerstep{2. Explain why changing two factors obscures the contrast.}
A score is $S(m,p,D)$ for model $m$, protocol $p$ and item set $D$. With fixed $D$,
\[
\begin{aligned}
S(m_2,p_2,D)-S(m_1,p_1,D)
={}&[S(m_2,p_2,D)-S(m_2,p_1,D)]\\
&+[S(m_2,p_1,D)-S(m_1,p_1,D)].
\end{aligned}
\]
The observed difference combines a protocol contrast and a version contrast. Measuring just the endpoints does not supply both terms. A matched version comparison uses the same protocol and items, or a design that separately estimates the factors.
\answerstep{3. Reverse a ranking without changing any answer.}
Suppose fixed answers have fact and clarity scores as follows:
\[
\begin{array}{c|cc}
&\text{Fact score}&\text{Clarity score}\\\hline
A&0.9&0.4\\B&0.7&0.9
\end{array}
\]
Define a combined score $S_w=wF+(1-w)C$, where $0\leq w\leq1$ is the weight on facts. At $w=0.8$, $S_A=0.80$ and $S_B=0.74$. At $w=0.2$, $S_A=0.50$ and $S_B=0.86$. The ranking reverses solely because the evaluation objective changes.
\answerstep{4. Interpret an uncertainty summary.}
For independent matched item differences $d_j$, estimate the mean difference by $\bar d$ and its SE by $s_d/\sqrt n$. For invented differences $(1,-1,2,0)$, $\bar d=0.5$, $s_d^2=5/3$ and SE $\sqrt{5/3}/2\approx0.645$. These are score units, not probabilities. Repeated generations or related items require an analysis consistent with their dependence.
\answerconnection{A presentation may change an encounter without improving its content. Indexing the model, protocol, items and score prevents that change from being mislabelled as a single model effect.}

\solutionslide{09}{Turn a model suggestion into a study}
\biomedicalanswer
\task{Draft a comprehension question about an uncertain scan finding. Describe a consequential misunderstanding and a population in which to check the wording.}
\answerstep{1. Ask what the explanation is meant to communicate.}
Suppose an explanation states that the scan contains an abnormality but its cause cannot be established from this image alone. A comprehension question could be: ``What does the scan tell us, and what does it leave unresolved?'' An adequate answer identifies the visible finding and distinguishes it from a confirmed cause. The question should be anchored to the actual explanation and source finding.

A consequential misunderstanding would be: ``The image proves that the abnormality is cancer,'' when the source has not established that diagnosis. Another would be: ``An uncertain cause means there is definitely no problem.'' Both replace uncertainty with unwarranted certainty in opposite directions.
\answerstep{2. Build a usable scoring guide.}
Specify in advance whether an answer must identify the finding, retain uncertainty and avoid a false diagnostic inference. Allow equivalent ordinary-language descriptions. Avoid giving credit simply for repeating a technical word. Ask a follow-up such as ``What makes you say that?'' to distinguish understanding from a guessed response.

Use independent assessors with examples of acceptable and unacceptable answers. If the research question concerns communication, score recipients' understanding as well as the factual accuracy of the text they received.
\answerstep{3. Check the question in the intended population.}
Test the wording with people resembling the intended recipients, including relevant language backgrounds, levels of health literacy and prior experience of imaging. Cognitive interviews can reveal whether words such as ``finding'' or ``uncertain'' mean what the researchers intend. Anxiety and prior knowledge may affect interpretation and should be considered in recruitment and measurement.
\answercheck{A useful study question specifies the misunderstanding it is designed to detect. Satisfaction, trust and factual comprehension can all matter, but they require distinct observations rather than one interchangeable score.}

\computationalanswer{09}
\task{Specify a causal outcome, construct an aggregate reversal across two strata and distinguish randomized participants from repeated question responses.}
\answerstep{1. Define the estimand before the data.}
Let $Y_i(1)$ be participant $i$'s comprehension score under an AI-assisted explanation and $Y_i(0)$ the score under the comparison. For example, define the score as the proportion of five predeclared questions answered correctly, ranging from zero to one. The target average effect is
\[
\Delta=\mathbb E[Y_i(1)-Y_i(0)].
\]
Only the assigned condition's outcome is observed for a participant. Random assignment supports estimating the average contrast under a well-defined intervention, no interference and the design's sampling and follow-up assumptions. It does not reveal the two potential outcomes for one person.
\answerstep{2. Construct the reversal explicitly.}
The following invented table records success on a single comprehension criterion. It illustrates the arithmetic of unequal group composition, not an estimated effect of AI:
\begin{center}
\begin{tabular}{lcc}
\toprule
Stratum & Group A & Group B\\\midrule
Novice & $60/100=60\%$ & $5/10=50\%$\\
Experienced & $9/10=90\%$ & $80/100=80\%$\\\bottomrule
\end{tabular}
\end{center}
Within each stratum A has the higher success rate. Aggregating gives
\[
P_A=\frac{60+9}{110}\approx0.627,\qquad
P_B=\frac{5+80}{110}\approx0.773.
\]
B now appears better because it contains many more experienced readers. The aggregate difference mixes within-stratum performance with composition.
\answerstep{3. Compare the same target mixture.}
If the target population is half novice and half experienced, standardizing both groups to those weights gives
\[
P_A^*=0.5(0.6)+0.5(0.9)=0.75,\qquad
P_B^*=0.5(0.5)+0.5(0.8)=0.65.
\]
This does not itself make observational rates causal. It shows why the population and weighting rule must be declared. Randomization balances composition in expectation, not necessarily exactly in a realized sample; stratification can address an important known factor by design.
\answerstep{4. Identify the analysis unit.}
If participants are assigned and each answers five questions, five responses are repeated measurements within one participant. Define a participant score or use an analysis accounting for this clustering. Counting five answers as five independently randomized people misstates the design. Freeze scoring and handling of missing answers before examining group results.
\answerconnection{The biomedical route defines a misunderstanding worth measuring. The formal route ensures that the measured comparison addresses the intended people and intervention rather than a difference in group composition.}

\solutionslide{10}{From patient data to a useful decision}
\biomedicalanswer
\task{Trace one input from acquisition to prediction to action, and explain how the care process could produce apparent accuracy.}
\answerstep{1. Trace an input and its availability.}
Imagine a research model predicting a later disease-change outcome from an MRI available at a specified decision date. The image is acquired, reconstructed, checked for quality, aligned or otherwise processed, and converted into input features. The model produces a forecast. A separate clinical process interprets that forecast with symptoms, previous findings, uncertainty and the burdens of possible actions.

The acquisition date is not always the availability date of every derived input. A report, annotation or later corrected record may become available only after the intended prediction time. A fair prospective task uses information actually available then.
\answerstep{2. Identify a care-process shortcut.}
Suppose more concerning cases are sent for an additional acquisition or an urgent review. A model might learn the associated protocol label or recording pattern. That feature could predict the outcome in the original service while partly reflecting clinicians' prior concerns and organizational decisions. If referral or recording practice changes, the association may weaken.

An even clearer failure is including a later note describing the outcome being predicted. The apparent forecast can then partly reconstruct a known result. High retrospective accuracy does not repair the time-order error.
\answerstep{3. Connect the forecast to a meaningful decision.}
Ask what an additional assessment would achieve, what burdens it entails and how errors would matter. A probability estimate does not choose those consequences. Evaluate the proposed use at the actual point in the workflow and compare it with the current procedure, including its existing clinical judgment.
\answercheck{A satisfactory trace identifies the physical measurement, processing, availability time, prediction target and decision. It gives a concrete route by which organizational information could substitute for the intended disease feature.}

\computationalanswer{10}
\task{Derive the action boundary for false-positive cost 1 and false-negative cost 4, and explain why a future note changes a prediction problem.}
\answerstep{1. Define the information set and probability.}
At time $t$, let $x_{\leq t}$ contain only available information. For future binary outcome $Y_{t+h}$ with horizon $h>0$, write
\[
p=P_\theta(Y_{t+h}=1\mid x_{\leq t}).
\]
The parameter $\theta$ indexes the fitted model. Interpreting its output as this probability requires evaluation in the relevant setting; algebra alone does not establish calibration.
\answerstep{2. Write the full loss table.}
For actions ``act'' and ``wait'', stipulate loss zero for the two correct decisions, $C_{FP}=1$ for acting without the event and $C_{FN}=4$ for waiting despite it. Expected losses are
\[
L_{\rm act}=P(Y=0\mid x)C_{FP}=(1-p),\qquad
L_{\rm wait}=P(Y=1\mid x)C_{FN}=4p.
\]
Both quantities use the same stipulated cost units.
\answerstep{3. Rearrange the comparison one step at a time.}
Acting has strictly lower expected loss when
\[
1-p<4p\quad\Longleftrightarrow\quad1<5p
\quad\Longleftrightarrow\quad p>\tfrac15=0.2.
\]
At $p=0.2$, both expected losses equal $0.8$. At $p=0.1$, acting costs $0.9$ in expectation and waiting $0.4$; at $p=0.3$, the values are $0.7$ and $1.2$. These substitutions check both sides of the boundary.

For arbitrary positive costs, the same rearrangement gives $p>C_{FP}/(C_{FP}+C_{FN})$. Scaling both costs by the same positive constant leaves this threshold unchanged. Nonzero correct-decision costs or additional actions require a revised comparison.
\answerstep{4. Explain information leakage.}
If a future note $N_{t+h}$ is included, the evaluated quantity is closer to $P(Y_{t+h}=1\mid x_{\leq t},N_{t+h})$. This is a different conditional distribution using information unavailable at the declared decision time. Renaming that input as an earlier feature does not make the prospective task valid.

\coursenotebook{15_ppv_and_decisions.ipynb}{Notebook 15} separates an event probability, a test-score classification threshold and a consequence-based action boundary. A forecast under observed care also does not by itself identify what would happen under a different treatment.
\answerconnection{Clinical reasoning supplies the relevant information date and consequences. The inequality makes the decision implied by those assumptions explicit and checkable.}

\solutionslide{11}{Keep the task fixed across Marr's levels}
\biomedicalanswer
\task{Explain why identical contours do not establish identical recognition, and what history-taking adds to the task.}
\answerstep{1. Separate the output from the process.}
Suppose a clinician and a model return the same outline. That is evidence about the contour on this case. It does not determine how either system selected the boundary. A model could exploit a reliable texture cue; the clinician could combine several contrasts, anatomical relationships and prior knowledge. The contour alone does not let us reconstruct either process.

Marr's levels ask about the purpose of a computation, the representations and procedures that carry it out, and its physical realization \citep{marr2010}. These are complementary questions about a fixed task. They should not be replaced by comparisons between different tasks simply because both involve images.
\answerstep{2. Explain what history changes.}
For lesion follow-up, history may identify why the scan was obtained, which symptoms matter, whether a relevant intervention occurred and how an earlier finding was interpreted. This context can change which comparison is clinically meaningful and how a boundary change is understood. It need not change the pixels of the current scan.

For example, the purpose may be documenting reproducible visible extent, rather than identifying every infiltrating cell. The clinician's context helps interpret whether that summary addresses the person's problem. An identical contour does not establish that the model incorporated that context.
\answerstep{3. Suggest evidence about recognition.}
Compare responses when image contrast, anatomical context or a suspected shortcut changes while the task stays fixed. For the clinician, ask which additional information is sought; for the model, test whether the relevant change alters the prediction. A spoken account is one observation and should be considered alongside behaviour, rather than treated as a complete readout of mechanism.
\answercheck{A sound answer distinguishes ``same contour'', ``same summary measurement'', ``same recognition procedure'' and ``same clinical purpose''. The first two can be checked directly without assuming the latter two.}

\computationalanswer{11}
\task{Derive the volume-disagreement bound and construct equal-volume masks at different locations.}
\answerstep{1. Define the masks and their physical units.}
On a common grid, let $m_i,\tilde m_i\in\{0,1\}$ indicate whether voxel $i$ belongs to either mask. Its physical volume is $v_i>0$ in mm$^3$. The selected volumes are
\[
V(m)=\sum_i v_i m_i,\qquad V(\tilde m)=\sum_i v_i\tilde m_i.
\]
Mask values are dimensionless; these sums are in mm$^3$. Dividing by 1000 converts either volume to mL.
\answerstep{2. Apply the triangle inequality.}
Subtract the sums, place the difference inside an absolute value and use that the absolute value of a sum cannot exceed the sum of absolute values:
\[
\begin{aligned}
|V(m)-V(\tilde m)|
&=\left|\sum_i v_i(m_i-\tilde m_i)\right|\\
&\leq\sum_i|v_i(m_i-\tilde m_i)|\\
&=\sum_i v_i|m_i-\tilde m_i|.
\end{aligned}
\]
The last equality uses $v_i>0$. For binary masks, $|m_i-\tilde m_i|$ is one exactly where the masks disagree. The right side is therefore their disagreement volume, also called the volume of their symmetric difference.
\answerstep{3. Show that equal volume can hide total displacement.}
Take two distinct voxels, each of volume 1~mm$^3$, and masks
\[
m=(1,0),\qquad\tilde m=(0,1).
\]
Both have volume 1~mm$^3$, so the left side of the bound is zero. They share no selected voxel, and the right side is $1+1=2$~mm$^3$. Their equal volumes do not establish spatial agreement. Missing volume and extra volume cancel in the signed total.

For a nested pair, $m=(1,1)$ and $\tilde m=(1,0)$, both sides instead equal 1~mm$^3$. The bound can therefore be exact; it is not a claim that all differences cancel.
\answerstep{4. State which information is discarded.}
The map $m\mapsto V(m)$ is not one-to-one: many spatial masks have the same summary. Consequently, equal outputs of the composed map $x\mapsto V(f_\theta(x))$ do not imply equal masks $f_\theta(x)$, still less equal algorithms or mechanisms. For comparisons on different grids, first define a physically consistent common representation; equal array indices do not automatically identify the same locations.

\coursenotebook{00_intro_timeline.ipynb}{Notebook 00} provides the explanatory-level context for these mappings.
\answerconnection{History determines which summary is meaningful for the task. The bound explains why agreement on that summary can leave substantial disagreement in the underlying spatial representation.}

\solutionslide{12}{Read about AI with an application in mind}
\biomedicalanswer
\task{Compare a proposed application with the present workflow. Describe one helpful change and one plausible failure.}
\answerstep{1. Define a concrete application and comparison.}
Consider assistance in drafting a longitudinal imaging summary from verified findings and prior reports. The current workflow is a professional reading the source material and composing the summary. The proposed workflow adds a model-generated draft, followed by review against the same sources. The comparison should concern an accepted summary, rather than the first text produced.

This choice makes the potential benefit specific: assistance might help organize dates and measurements consistently, leaving more attention for interpretation and communication. Whether it actually saves time or improves clarity must be observed. A broad claim that ``AI improves imaging'' is too imprecise to assess.
\answerstep{2. Explain a plausible failure.}
A draft might convert ``no definite increase'' into ``increase'', transpose the side of a finding or merge measurements from different dates. The resulting prose may read smoothly while changing clinical meaning. A second possibility is that the tool works well on common report formats but fails on an unusual sequence or phrasing.

The user must therefore be able to trace statements to their source. An unsupported sentence should remain visible as an error even if the overall summary is impressive.
\answerstep{3. Identify an informative assessment.}
Use matched source cases and compare accepted summaries for factual errors, preservation of uncertainty, reader comprehension and total work time. Record errors by type and important case group, not just their overall average. Include representative difficult cases rather than selecting only successful demonstrations.

The application discussion in the syllabus motivates coexistence with professional work \citep{mollick2026coexistence}. It does not supply performance estimates for this particular proposed use. The worked answer is a design for evaluating that use.
\answercheck{A good answer names the baseline workflow, the exact change assistance introduces, a failure with a clear consequence and an observation that could reveal both benefit and failure.}

\computationalanswer{12}
\task{Reweight group loss differences for another population and specify a controlled intervention on a suspected irrelevant feature.}
\answerstep{1. State the sign convention.}
For matched case $i$, define $d_i=L_i^{AI}-L_i^B$, where $B$ labels the baseline. Loss is smaller when performance is better, so negative $d_i$ favours AI. Suppose two groups have fixed mean differences $\bar d_1=-0.1$ and $\bar d_2=+0.5$ in a common loss scale.
\answerstep{2. Change only the population weights.}
Let $w$ be the target population's proportion in group 1; group 2 has proportion $1-w$. The aggregate mean is
\[
\bar d(w)=w(-0.1)+(1-w)(0.5)=0.5-0.6w.
\]
For the syllabus's mixture $w=0.9$, this is $-0.04$. For equal proportions, it is $+0.20$. For $w=0.6$, it is $+0.14$. The same two group-specific contrasts can thus support an aggregate improvement or deterioration in different populations.

Solve the sign boundary:
\[
0.5-0.6w<0\quad\Longleftrightarrow\quad w>\frac56.
\]
At $w=5/6$ the aggregate ties. This calculation assumes the two group-specific means remain applicable in the new setting; a new population can also differ within groups. Weights alone cannot correct that additional shift.
\answerstep{3. Define the feature intervention.}
In a constructed marker experiment, write input $x=(b,m)$ for biological content $b$ and marker $m$. Keep $b$, its target $y$ and the fixed predictor $f$ unchanged. Compare $f(b,m)$ with $f(b,m')$ after changing only the marker. Record how often the label or score changes and how that affects the stipulated loss. \coursenotebook{08_attention_maps.ipynb}{Notebook 08} implements a marker-intensity replacement while preserving the central pattern in a transparent parameterized model. \coursenotebook{11_shortcut_dependence.ipynb}{Notebook 11} instead changes marker--target association in an analytic binary construction.

If an altered marker changes the output while the intended target stays fixed, the experiment establishes sensitivity to that marker under the intervention. It does not by itself quantify clinical harm. In real images, check that the change does not corrupt the disease information or create an unrealistic artefact; otherwise interpretation becomes ambiguous.
\answerstep{4. Distinguish the two questions.}
Reweighting asks how fixed group contrasts combine for a declared target population. The feature intervention asks what the model's function depends on. One is a statement about averaging; the other is a statement about responses to a controlled input change. Both can be relevant to evaluating the same application.
\answerconnection{Clinical examples identify which errors and people matter. The formal decomposition shows how an apparently favourable average can coexist with a failure in one group or dependence on an unintended cue.}

\solutionslide{13}{Preserve the ability to judge}
\biomedicalanswer
\task{Explain how repeated acceptance of an overlay could affect later unaided judgments, and propose an activity requiring explanation and checking.}
\answerstep{1. Describe two possible learning mechanisms.}
An overlay can provide feedback: the learner compares a prediction with a proposed boundary, notices a missed feature and revises an explanation. It can also remove practice if the learner accepts it before inspecting the image. In the second workflow, fewer opportunities remain to decide where an uncertain boundary lies or to recognize a plausible-looking error.

These mechanisms motivate different predictions about later unaided performance. Assistance could help, harm or leave it unchanged depending on how it is used and what the learner practices. The mathematics-learning study cited in the syllabus is a reason to ask this question, not an effect estimate for imaging trainees \citep{bastani2025learning}.
\answerstep{2. Design a teaching activity.}
First show the image without an overlay. Ask the learner to mark a boundary, state uncertainty and name the contrast or anatomical relationship supporting the choice. Then reveal the proposed overlay and ask the learner to explain two agreements and one possible disagreement. The learner must check the disagreement against relevant views or source information before accepting a final result.

Include some deliberately imperfect proposals in a declared teaching exercise so that checking is meaningful. Discuss why a correction is justified, rather than rewarding conformity with the proposed contour.
\answerstep{3. Assess retained judgment independently.}
Later, present new cases without overlays or access to earlier assisted answers. Use a fixed rubric that addresses boundary quality and reasoning, including uncertainty. Compare groups with matched teaching time and baseline assessment where feasible. The later assessment should measure the skill of interest, not recall of a previously corrected contour.
\answercheck{The answer should explain what practice assistance changes, require an unaided first judgment and identify a later assessment that can reveal retained competence. Immediate speed and later learning are separate outcomes.}

\computationalanswer{13}
\task{Distinguish percentage points from relative percent change and design a participant-level assessment of later unaided performance.}
\answerstep{1. Calculate the two outcome contrasts.}
For the syllabus's invented means, assisted performance is $0.80$ in group A and $0.60$ in group B, while later unaided performance is $0.60$ and $0.75$. Therefore,
\[
\Delta_a=0.80-0.60=+0.20,\qquad
\Delta_l=0.60-0.75=-0.15.
\]
The positive assisted contrast and negative later contrast are logically compatible: they concern different outcomes. Neither is obtained by subtracting a participant's assisted score from a different person's later score.
\answerstep{2. Translate the numerical scales carefully.}
For proportions correct, $-0.15$ means a decrease of 15 percentage points, from 75\% to 60\%. Relative to the comparison mean 0.75, the change is
\[
100\frac{0.60-0.75}{0.75}=-20\%.
\]
The denominator determines the relative percentage. The assisted contrast of $+0.20$ is a 20-percentage-point increase, or a relative increase of $100(0.80-0.60)/0.60\approx33.3\%$.

These constructed numbers are not the cited study's reported 17\% reduction. Keeping them explicitly labelled prevents a teaching calculation from being mistaken for a published finding.
\answerstep{3. Define the assessment and estimator.}
For participant $i$ in assigned group $g\in\{A,B\}$, let $S_{l,ig}$ be the proportion correct on a new, unaided assessment using a fixed scoring rule. If $n_g$ participants are observed, estimate retained-performance contrast by
\[
\widehat\Delta_l=
\frac1{n_A}\sum_{i=1}^{n_A}S_{l,iA}
-\frac1{n_B}\sum_{i=1}^{n_B}S_{l,iB}.
\]
The unit is a participant, even if each completes many images or questions. An analysis of individual responses must represent their dependence within participants, and possibly shared cases as well.
\answerstep{4. State what the design must preserve.}
Use new items, remove access to assisted answers, keep assessment conditions and scoring comparable, and blind scoring to group where possible. Record baseline skill and missing later assessments. Random assignment and appropriate follow-up support a causal interpretation of the teaching contrast; an observational comparison needs additional assumptions about group differences.

Without sample sizes and participant variation, the two means alone do not supply an uncertainty interval. A precise-looking percentage is not evidence of a precisely estimated effect.
\answerconnection{The biomedical hypothesis concerns what the learner continues to practice. Separate outcomes and participant-level comparisons are needed to test whether that practice changes retained judgment.}

\solutionslide{14}{Map the frontier task by task}
\biomedicalanswer
\task{Describe two apparently similar tasks requiring different clinical knowledge, and a setting in which success on one offers weak reassurance about the other.}
\answerstep{1. Specify the two questions.}
Task A is outlining a conspicuous abnormal region on an image according to an agreed annotation policy. Task B is deciding why a region has changed between scans after an intervention. Both involve inspecting a lesion, but the second question requires interpretation of time, prior treatment, image contrasts and other clinical observations. A good outline does not answer the causal question.

Another valid pair is recognizing a common large finding and detecting a small abnormality in an unusual location. Their apparent visual similarity can conceal different spatial scales, error consequences and background anatomy.
\answerstep{2. Predict a revealing circumstance.}
A model evaluated mainly on familiar acquisitions and untreated cases may encounter postoperative anatomy, a changed acquisition protocol or an atypical appearance. An outline could remain technically plausible while its interpretation is wrong for that context. The professional question is then what the changed appearance means, not simply whether a contour can be drawn.

This is a concrete version of the uneven capability frontier discussed in the syllabus \citep{mollick2026coexistence}. It concerns the task and setting, rather than a single general label such as ``good at medical images''.
\answerstep{3. Name a useful evaluation.}
Describe each task's target, relevant population and error type separately. Examine performance under the specific changed circumstances, including what information was available and how the reference conclusion was established. If the issue is interpretation of change, add observations that distinguish competing causes rather than relying only on overlap with a contour.
\answercheck{The strongest answers identify the extra knowledge needed for the second task and explain why the first task's evaluation did not test it. A vague statement that ``hard cases are different'' leaves the transfer question unspecified.}

\computationalanswer{14}
\task{For a fixed marker-based rule, derive accuracy and risk under changing marker--target association. Identify exactly what stays fixed.}
\answerstep{1. Define the rule and loss.}
Let input $X$ contain a binary marker $M\in\{0,1\}$ and other features. The target is $Y\in\{0,1\}$. Fix the predictor
\[
f(X)=M,\qquad L(f(X),Y)=\mathbf1[M\ne Y].
\]
The function copies the marker irrespective of other content. Its coefficients, threshold and procedure do not change between settings.
\answerstep{2. Convert the expectation into a probability.}
Since an indicator's expectation equals the probability of its event,
\[
\begin{aligned}
\operatorname{Accuracy}_P(f)
&=\mathbb E_P\mathbf1[M=Y]=P(M=Y),\\
R_P(f)&=\mathbb E_P\mathbf1[M\ne Y]=P(M\ne Y)=1-P(M=Y).
\end{aligned}
\]
The last equality holds because equality and inequality are complementary events for these labels.
\answerstep{3. Give fully specified changed populations.}
For a constructed family, keep $P(Y=0)=P(Y=1)=1/2$. Given either target value, let the marker match it with probability $q$, and take the other binary value with probability $1-q$. This specifies a valid joint distribution with
\[
P(M=Y)=q,\qquad \operatorname{Accuracy}=q,\qquad R=1-q.
\]
At $q=0.95$, accuracy is $0.95$ and risk $0.05$. At $q=0.5$, both are $0.5$. At $q=0.05$, accuracy is $0.05$ and risk $0.95$. Class prevalence stays balanced throughout, so the changed marker association alone suffices for the reversal.
\answerstep{4. Relate this to distribution shift.}
The changed object is the joint law $P(M,Y)$, equivalently the association of marker and target. The predictor is fixed. The numerical example therefore separates a change in setting from retraining or an architectural change. \coursenotebook{11_shortcut_dependence.ipynb}{Notebook 11} provides the corresponding constructed experiment.

The extreme example in the syllabus also makes the point: a rule correct only at input $a$ has risk zero when the population is entirely $a$, and risk one when it is entirely input $b$, which the rule labels incorrectly. Claims about transfer need assumptions connecting the source and target populations or direct target evaluation.
\answerconnection{Clinical context identifies why a second task or setting may differ. The fixed-rule calculation demonstrates how performance can change substantially without any change to the model itself.}

\solutionslide{15}{Specify, inspect, revise and stop deliberately}
\biomedicalanswer
\task{Rewrite an uncertain finding for a patient. Explain which wording preserves uncertainty and which additional observation could justify a stronger statement.}
\answerstep{1. Start from a source that actually contains uncertainty.}
For this hypothetical exercise, suppose the source says: ``A small focus of altered signal is present. Its cause is nonspecific on this examination.'' A faithful plain-language version is: ``The scan shows a small area that looks different from the surrounding tissue. This scan alone does not tell us what caused it. Its meaning needs to be considered with the rest of the information about you.''

The rewrite identifies the observation without inventing a diagnosis, a symptom link or a management instruction. Different wording is acceptable if it preserves those distinctions.
\answerstep{2. Explain why the uncertainty words matter.}
``Shows'' applies to the visible appearance. ``This scan alone'' limits the evidential claim. ``Does not tell us what caused it'' prevents an abnormal image appearance from becoming a confirmed disease label. Removing these phrases can alter the proposition rather than merely improve the style.

The answer should also remain understandable. Replacing ``nonspecific'' with another unexplained technical word does not solve the communication problem.
\answerstep{3. Name evidence relevant to a stronger claim.}
An earlier comparable scan, an additional appropriate measurement or clinical information might help distinguish competing explanations. Specify which alternative that observation would address. For example, a reliable longitudinal comparison could establish that the visible feature changed, but change alone would not necessarily establish its cause. Stronger wording must match what the new evidence actually resolves.

Before accepting the explanation, compare it sentence by sentence with the source: finding, location, chronology, negation and uncertainty. Ask a recipient to explain what they think is known and unknown.
\answercheck{A correct rewrite can be simple without being more certain. Repeated polishing is useful only if the supported meaning survives each revision.}

\computationalanswer{15}
\task{Construct a grammatically acceptable but unsupported statement, strengthen its acceptance predicate, and explain why repeated deterministic checks add no new evidence.}
\answerstep{1. Give a counterexample to a weak check.}
Let state $q$ contain the source finding, draft explanation and check results. Suppose $C_g(q)=1$ means that grammar and sentence structure are acceptable. The draft ``This scan confirms cancer'' is grammatical, so it may satisfy $C_g=1$ while contradicting the nonspecific source used above. Thus
\[
C_g(q)=1\quad\not\Rightarrow\quad\text{the draft is supported by the source}.
\]
This is a logical counterexample, not a measured error rate of a language tool.
\answerstep{2. Define a more relevant predicate.}
Let $C_e(q)$ check whether the stated findings, location, chronology, negation and certainty agree with the source. Let $C_s(q)$ check whether factual claims have identifiable supporting source passages. Define
\[
C(q)=C_g(q)\land C_e(q)\land C_s(q).
\]
The symbol $\land$ means logical ``and'': every component must return true. For the unsupported diagnosis, a correctly applied $C_e$ returns zero. These checks need an operational rubric; merely naming a function ``evidence check'' does not establish that its implementation detects every unsupported claim.
\answerstep{3. Work through the stopping rule.}
For budget $K=2$, suppose $q_0$ contains the unsupported claim, $q_1$ a more elegant version of the same claim and $q_2$ the faithful uncertain explanation. If the respective acceptance results are $0,0,1$, then
\[
\tau=\min\{k\in\{0,1,2\}:C(q_k)=1\}=2.
\]
If all three results are zero, the set is empty and the procedure reports no accepted result within the budget. It must not silently return the final draft as if it had passed. Choosing an extended-real convention $\tau=+\infty$ is possible, but must be declared; it is not the finite minimum above.
\answerstep{4. Explain the evidential limit of repetition.}
For a deterministic check on unchanged state $q$, each repetition returns the same value $C(q)$. Repeating it $N$ times does not create $N$ independent observations of correctness or test a new property. Changing the draft, evidence or check can produce a new test, but its contribution should be stated explicitly.

Acceptance establishes compliance with the selected, implemented criteria. It does not establish every factual, biological or communicative property that might matter.
\answerconnection{The biomedical answer supplies what must be preserved. The formal answer translates those requirements into an inspectable procedure and exposes the limits of a weak acceptance rule.}

\solutionslide{16}{Compartments, anatomy and measurement}
\biomedicalanswer
\task{Predict how a closed fluid barrier changes spread and explain how coarse sampling can change the pattern without changing growth rates.}
\answerstep{1. State what ``closed'' means in this model.}
The lesson assigns some fluid-labelled regions as barriers to cell transport. Under that assumption, cells cannot cross a closed interface, although proliferation can still increase their number on the occupied side. The population may spread along connected tissue routes around the barrier instead of taking the shortest straight path through it.

If a barrier completely separates an unseeded component from every seeded component, and the model has no other cell source there, that unseeded component stays empty. A partial barrier can instead delay or redirect access. These are predictions about the declared transport geometry, not universal statements about all tumour--fluid interactions.
\answerstep{2. Explain the effect of a coarse grid.}
A thin fluid layer may disappear or become a connected tissue voxel when anatomy is resampled. The discretized model now contains a path that the finer representation blocked. Cells can enter the newly connected region and then proliferate there under the same assigned growth rates. The altered pattern therefore need not come from altered cell biology; it can come from altered connectivity.

The reverse is also possible: coarse sampling can close a narrow tissue passage and impede spread. Grid spacing changes physical representation as well as the number of numerical unknowns.
\answerstep{3. Choose observations that distinguish the effects.}
Inspect the resampled tissue mask and the relevant physical slices, trace whether the pathway is open and compare simulations with matched initial fields and rates as closely as the grids permit. Report voxel dimensions and the barrier's physical thickness. If the goal is a pure solver-resolution test, use a geometry defined consistently across grids rather than silently changing the anatomical domain.

BrainWeb supplies a normal anatomical phantom; density, seed and rates are assigned model quantities in this lesson. Its image background is not observed tumour progression \citep{lundervold2026glioma}.
\answercheck{Explain the chain from anatomical sampling to connectivity, transport and local growth. A mass-conserving computation can still represent the wrong anatomical pathway.}

\computationalanswer{16}
\task{Prove that increasing a density cutoff cannot increase selected volume. Check the geometry and units of the density integral, threshold volume and cell count.}
\answerstep{1. Recover the physical voxel volume.}
An affine maps a homogeneous voxel coordinate to a physical coordinate:
\[
\begin{bmatrix}\mathbf x\\1\end{bmatrix}
=A\begin{bmatrix}\mathbf i\\1\end{bmatrix}.
\]
The spatial block $A_{3\times3}$ maps index displacements to millimetres. Its determinant is signed volume scaling; taking its absolute value gives $v=|\det A_{3\times3}|$ in mm$^3$. Translation changes positions, not voxel volume. For spacings 2, 3 and 4~mm on orthogonal axes, $v=24$~mm$^3$; reversing an axis changes the determinant's sign but not this positive volume.
\answerstep{2. Define the selected sets and prove inclusion.}
Let $u_i=c_i/K$ be dimensionless density fraction and $S_\tau=\{i:u_i\geq\tau\}$. If $\tau_2\geq\tau_1$, every $i\in S_{\tau_2}$ satisfies $u_i\geq\tau_2\geq\tau_1$, hence $i\in S_{\tau_1}$. Therefore $S_{\tau_2}\subseteq S_{\tau_1}$ and
\[
V_{\tau_2}=\frac v{1000}|S_{\tau_2}|
\leq\frac v{1000}|S_{\tau_1}|=V_{\tau_1}.
\]
Here $|S|$ counts elements; $v/1000>0$ converts each voxel's volume to mL. With unequal positive voxel volumes, the same argument uses sums over the nested sets. The proof requires a fixed field and geometry, not a rerun after changing model parameters.
\answerstep{3. Compute three different summaries.}
For four 4-mm isotropic voxels, $v=64$~mm$^3=0.064$~mL. Take $u=(0.10,0.25,0.60,1.00)$, whose sum is $1.95$. Then
\[
\begin{aligned}
B&=0.064(1.95)=0.1248\ \mathrm{mL},\\
V_{0.2}&=0.064(3)=0.192\ \mathrm{mL},\\
V_{0.5}&=0.064(2)=0.128\ \mathrm{mL}.
\end{aligned}
\]
The density integral $B$ weights each voxel by its fraction; the threshold volumes count every selected voxel fully. For constant $K$ in cells/mm$^3$, the cell count is
\[
N=Kv\sum_i u_i=K\,(124.8\ \mathrm{mm}^3).
\]
Its units are cells, not mL. Multiplying the displayed mL value $B$ by $K$ without the factor 1000 would give the wrong unit conversion.
\answerstep{4. Interpret the invariant.}
Changing only $\tau$ can change $V_\tau$ while leaving $B$ and $N$ unchanged. The notation $B$ here means a density integral; on slide 19 it means a local birth rate. \coursenotebook{14b_brainweb_growth.ipynb}{Notebook 14b} keeps geometry, assigned parameters and simulated fields distinct.
\answerconnection{The biological answer explains how a barrier shapes a pattern. The calculations distinguish changes in that pattern from changes caused only by threshold or physical-volume accounting.}

\solutionslide{17}{Population growth and a mechanistic hypothesis}
\biomedicalanswer
\task{Predict the outcomes of disabling movement and disabling proliferation in a closed domain. Identify an observation separating accumulation from redistribution.}
\answerstep{1. Disable movement while keeping proliferation.}
In a model without transport, cells cannot seed a previously empty location by moving into it. Where cells are initially present, their density can rise toward the local limiting scale. An initially sparse tail may become visible as its existing cells proliferate; that appearance is not proof that new cells migrated there. An exactly empty location remains empty in this deterministic model without another source.

Thus the expected picture is local accumulation at occupied sites, with no transport across their boundaries. Whether a threshold contour expands depends partly on the initial field, including low-density tails.
\answerstep{2. Disable proliferation while keeping movement.}
With a closed, no-flux boundary, movement redistributes the existing population. High-density regions can become less dense while neighbouring regions gain cells. The total cell number remains constant, even though the apparent footprint can change. In disconnected compartments, each separately closed compartment retains its own total.

A thresholded display can grow or shrink during redistribution, so visible area alone is not a count of newly created cells. With both movement and proliferation disabled, the field remains unchanged.
\answerstep{3. Name a discriminating observation.}
Within the simulation, compare the spatial density field with its full-domain integral. Rising density at one site alongside falling density elsewhere and a constant total supports redistribution. An increasing total requires a source such as proliferation under this model's assumptions. In biological work, independent information about division and movement would be needed to separate those mechanisms; MRI appearance alone is not a direct cell count.

The go-or-grow hypothesis additionally distinguishes migratory and proliferative states \citep{hatzikirou2012}. Basic Fisher--KPP combines effective dispersal and logistic growth without an explicit switching variable.
\answercheck{Use local observations and a total together. An enlarging coloured region does not by itself distinguish migration, proliferation or a display-threshold effect.}

\computationalanswer{17}
\task{Derive the no-flux total-cell balance, establish conservation without proliferation and verify the diffusion-free logistic solution.}
\answerstep{1. Define the field and check units.}
On fixed tissue domain $\Omega$, let $c(\mathbf x,t)$ and constant capacity $K>0$ be in cells/mm$^3$, with $\mathbf x$ in mm and $t$ in days. Diffusivity $D\geq0$ has units mm$^2$/day and rate $\rho\geq0$ has units day$^{-1}$. Fisher--KPP is
\[
\partial_t c=\nabla\!\cdot(D\nabla c)+\rho c(1-c/K).
\]
The gradient has units cells/mm$^4$. Flux $\mathbf j=-D\nabla c$ therefore has units cells/mm$^2$/day; its divergence, like each other term in the equation, has units cells/mm$^3$/day.
\answerstep{2. Integrate and apply the boundary condition.}
Define total $M(t)=\int_\Omega c\,dV$. For sufficiently regular fields on a fixed domain, differentiate under the integral and use the divergence theorem:
\[
\begin{aligned}
\dot M
&=\int_\Omega\nabla\!\cdot(D\nabla c)\,dV
+\int_\Omega\rho c(1-c/K)\,dV\\
&=\int_{\partial\Omega}\mathbf n\cdot D\nabla c\,dS
+\int_\Omega\rho c(1-c/K)\,dV\\
&=\int_\Omega\rho c(1-c/K)\,dV.
\end{aligned}
\]
Here $\partial\Omega$ is the boundary, $\mathbf n$ its outward unit normal and $dS$ a surface element. No flux means $\mathbf n\cdot\mathbf j=0$, equivalently $\mathbf n\cdot D\nabla c=0$. The total $M$ has units cells and $\dot M$ cells/day. Setting $\rho=0$ gives $\dot M=0$, hence $M(t)=M(0)$.
\answerstep{3. Verify the closed-compartment solution directly.}
With $D=0$, each location obeys $dc/dt=\rho c(1-c/K)$. For initial value $c_0>0$, set $a=K/c_0-1$ and $q(t)=a e^{-\rho t}$. The proposed solution is $c(t)=K/(1+q(t))$. Since $q'=-\rho q$,
\[
c'(t)=\frac{K\rho q}{(1+q)^2},\qquad
\rho c(1-c/K)=\rho\frac K{1+q}\frac q{1+q}
=\frac{K\rho q}{(1+q)^2}.
\]
The differential equation is satisfied, and
\[
c(0)=\frac K{1+(K/c_0-1)}=c_0.
\]
For $0<c_0<K$ and $\rho>0$, the solution increases toward $K$; $c_0=K$ gives a constant. The separate initial condition $c_0=0$ has solution zero, although the displayed formula contains $K/c_0$ and cannot be evaluated there directly.
\answerstep{4. Turn the results into verification checks.}
\coursenotebook{13_verified_growth_model.ipynb}{Notebook 13} can compare a transport-only run with constant total and a reaction-only run with this exact solution. Check errors at matched times and units. A numerical method can have finite discretization error while still approximating the continuum equations; these checks test implementation, not physiological adequacy.
\answerconnection{The qualitative predictions correspond to two separately testable mathematical limits: conserved redistribution and local logistic accumulation.}

\solutionslide{18}{An interactive simulation is an experiment}
\biomedicalanswer
\task{Predict where oxygen becomes scarce, change uptake alone, interpret viable and necrotic patterns, and distinguish a new run from a new view.}
\answerstep{1. State a spatial prediction and its assumptions.}
In the lesson's coupled model, viable cells consume oxygen and accumulated material suppresses an assigned delivery factor. A plausible prediction is that scarcity will develop in regions of high burden with limited replenishment, often within the denser interior of a growing population. The location need not be the geometric centre: tissue connectivity, initial seed and the chosen parameters can alter the pattern.

Mildly reduced oxygen can slow proliferation before the lower death threshold is reached. Greater scarcity can then transfer viable cells into the necrotic-material compartment. These are separate mechanisms, so a smaller viable population does not imply equally increased necrosis.
\answerstep{2. Describe a controlled notebook comparison.}
In \coursenotebook{14d_coupled_tumour_growth.ipynb}{Notebook 14d}, start with the small CPU preset and record its seed, domain, grid, rates, capacity and time settings. Compare uptake zero with uptake four while keeping the remaining settings fixed. These are instructions for an experiment, not new measurements claimed by this document.

Before running, predict the direction of change in oxygen, viable growth and necrotic accumulation. Afterward, inspect matched days and physical slices, and record full-domain trajectories as well as images. If oxygen never crosses the death threshold, reduced growth without much necrosis is compatible with the equations. An unexpected pattern is a reason to inspect assumptions and settings rather than reinterpret every dark region as dead tissue.
\answerstep{3. Explain a display change.}
After a run, move to another saved time or physical slice. This reveals another part of the computed trajectory; it does not change the parameters that generated it. Changing only a density display threshold changes what is visible, not the underlying full-volume cell count. The oxygen view uses the tissue mask rather than the tumour-density cutoff.
\answercheck{Record whether each control changes the model, the saved time being inspected or the display. Their effects answer different questions even when the screen changes in each case.}

\computationalanswer{18}
\task{Define a matched difference in full-volume viable-cell counts for uptake zero and four. Explain its invariance to display threshold.}
\answerstep{1. Separate the solution from its observation.}
Let $\theta_\beta$ denote the parameter vector whose uptake component is $\beta$. The solution operator gives
\[
U_{\theta_\beta}(t)=S_{\theta_\beta}(t;U_0)
=(c_\beta(\cdot,t),z_\beta(\cdot,t),o_\beta(\cdot,t)).
\]
The dot marks the full spatial argument. A displayed image is $I=H_{\tau,k,j}[U_{\theta_\beta}(t)]$, where $\tau$ is a cutoff, $k$ a slice and $j$ a selected field. The solution and the display operator are different mathematical objects.
\answerstep{2. Define a total with units.}
At one fixed physical time $t_*$ and on the same domain $\Omega$, define
\[
N_\beta(t_*)=\int_\Omega c_\beta(\mathbf x,t_*)\,dV,
\qquad \Delta N(t_*)=N_0(t_*)-N_4(t_*).
\]
For viable densities in cells/mm$^3$, these quantities are in cells. On equal-volume voxels, $N_\beta\approx v\sum_{i\in\Omega}c_{\beta,i}$. If the stored array is instead fraction $u=c/K$, multiply by the relevant $K$ before integrating. A sum of $u$ alone is not a cell count.

Positive $\Delta N$ means fewer viable cells in the uptake-four run at that day. It does not state that $\Delta N$ cells died, nor establish a pointwise ordering of every voxel. Its observed sign and size should be reported rather than inferred from a differently scaled picture.
\answerstep{3. Prove the display invariance.}
The expression for $N_\beta$ contains $c_\beta$, $\Omega$ and physical volume, but not $\tau$. Holding those objects fixed while replacing $H_\tau$ by $H_{\tau'}$ leaves $N_\beta$ and therefore $\Delta N$ unchanged. In contrast,
\[
V_{\tau,\beta}=\frac v{1000}\sum_i\mathbf1[c_{\beta,i}/K\geq\tau]
\]
does depend on the threshold. A total computed only from displayed voxels is this different, threshold-dependent measurement, not the full-domain count.
\answerstep{4. Specify what must be matched.}
Use the same initial viable and necrotic fields, tissue mask, spatial coordinates, $D$, $\rho$, $K$, other coupling settings, output days and declared solver settings. Match views independently when comparing images. Changing the grid also resamples anatomy and can change cell connectivity, so it is not by itself an isolated display or time-step intervention.
\answerconnection{The biological route predicts a mechanism-dependent pattern. The formal route defines a measurement that can test that prediction while remaining independent of how the pattern is displayed.}

\solutionslide{19}{From a biological diagram to coupled variables}
\biomedicalanswer
\task{Predict a viable rim and necrotic core, explain why missing viable cells need not equal necrosis, and distinguish both compartments from MRI signal.}
\answerstep{1. Explain a possible rim--core pattern.}
Where delivery can keep pace with consumption, viable cells can remain oxygenated and proliferate. In a region of high consumption and suppressed delivery, oxygen can fall far enough to slow births, and then below a lower threshold that activates death transfer. Under suitable settings this can leave a more viable outer region and accumulated necrotic material deeper within the population.

This is a conditional explanation, not a universal geometry. The lesson does not explicitly model vessels, angiogenesis or every form of tissue injury represented in a biological diagram \citep{comba2021heterogeneity}. Its assigned oxygen and material fields should therefore be read as model compartments.
\answerstep{2. Use a material-accounting example.}
Imagine two closed model populations starting with 100 viable cell equivalents and no necrotic material. Over an interval, the reference produces 40 births. The coupled population produces 10 births, transfers 15 viable equivalents to necrotic material and clears 5 necrotic equivalents. At the end:
\begin{itemize}
\item Reference: 140 viable equivalents.
\item Coupled: $100+10-15=95$ viable and $15-5=10$ necrotic equivalents.
\end{itemize}
The viable deficit is $140-95=45$, while remaining necrotic material is only 10. Thirty of the deficit are births that did not occur; fifteen reflect death transfer, five of which have already been cleared. These are invented accounting values, not notebook output.
\answerstep{3. Keep imaging interpretation separate.}
Viable density, necrotic material and oxygen are separate assigned fields. An MRI signal depends on acquisition and tissue properties; it is not one of these densities read out directly. A dark or coloured region in an overlay must be interpreted using its field label, units, threshold and legend. Calling it an MRI-measured necrotic volume would add an observation model that the lesson has not established.
\answercheck{A useful explanation distinguishes reduced births, transfer and clearance. It avoids reading a deficit relative to another simulation as a one-for-one map of cell death.}

\computationalanswer{19}
\task{Derive the total viable-plus-necrotic balance and explain exactly how it constrains interpretation of the viable deficit.}
\answerstep{1. State compartments, rates and boundary assumptions.}
On fixed tissue domain $\Omega$, let $c,z$ be viable and necrotic densities in cell equivalents/mm$^3$. For relative oxygen $o$, define local birth and transfer rates
\[
B=\rho g(o)c[1-(c+z)/K]_+,\qquad T=\delta h(o)c.
\]
Here $K>0$ is the common density scale, $D$ is diffusivity in mm$^2$/day, and $\rho,\delta\geq0$ are proliferation and death-transfer rates in day$^{-1}$. The dimensionless modifiers are $g(o)=\min(1,o/\theta_p)$ and $h(o)=\max(1-o/\theta_d,0)$, with oxygen thresholds $0<\theta_d<\theta_p<1$.
The positive part $[a]_+=\max(a,0)$ prevents negative births. Both $B$ and $T$ have units cell equivalents/mm$^3$/day. The governing density equations are
\[
\partial_tc=\nabla\!\cdot(D\nabla c)+B-T,\qquad
\partial_tz=T-\lambda_z z.
\]
Assume constant clearance rate $\lambda_z\geq0$ in day$^{-1}$, no transport of $z$, no cell flux across the boundary and enough regularity for the integral operations below. The shipped presets use $\lambda_z=0$; examples with clearance require a positive rate.
\answerstep{2. Add first, then integrate.}
At each position, adding gives
\[
\partial_t(c+z)=\nabla\!\cdot(D\nabla c)+B-\lambda_z z.
\]
The term $-T+T$ cancels because transfer changes compartment, not the amount represented in their sum. With $M(t)=\int_\Omega(c+z)\,dV$, the divergence theorem and no-flux condition give
\[
\begin{aligned}
\dot M
&=\int_{\partial\Omega}\mathbf n\cdot D\nabla c\,dS
+\int_\Omega B\,dV-\lambda_z\int_\Omega z\,dV\\
&=\int_\Omega B\,dV-\lambda_z\int_\Omega z\,dV.
\end{aligned}
\]
Each side has units cell equivalents/day. Births add material and clearance removes it. Transfer can affect later births through changed state, but it is not a direct loss term from $M$.
\answerstep{3. Distinguish a stock from accumulated transfer.}
Write $C=\int_\Omega c\,dV$, $Z=\int_\Omega z\,dV$, $\mathcal B=\int_\Omega B\,dV$ and $\mathcal T=\int_\Omega T\,dV$. Then
\[
\dot C=\mathcal B-\mathcal T,\qquad
\dot Z=\mathcal T-\lambda_z Z.
\]
For $Z(0)=0$, integrating the second equation gives
\[
Z(t)=\int_0^t e^{-\lambda_z(t-s)}\mathcal T(s)\,ds.
\]
The exponential discounts material cleared since time $s$. With $\lambda_z=0$, $Z(t)$ equals cumulative transfer under this zero-initial-necrosis assumption. With clearance it is generally smaller.
\answercontinued{19}
\answerstep{4. Relate the reference deficit to its components.}
For a matched no-death reference $C_r$ with the same initial viable total, $\dot C_r=\mathcal B_r$. Subtract and integrate:
\[
C_r(t)-C(t)
=\int_0^t(\mathcal B_r-\mathcal B)\,ds
+Z(t)+\int_0^t\lambda_z Z(s)\,ds.
\]
The deficit combines a difference in cumulative births, remaining necrotic material and cleared material. The birth difference can reflect oxygen effects and the changed density-dependent crowding. No pointwise death map follows from a single viable-difference image. If initial totals differ or material crosses the boundary, additional terms are required.

\answerstep{5. Check a simple transfer-only limit.}
For a closed uniform compartment, disable births and clearance and stipulate a constant effective transfer rate $q=0.2$ day$^{-1}$. The isolated transfer equations have solutions
\[
C(t)=100e^{-qt},\qquad Z(t)=100(1-e^{-qt})
\]
for initial totals $C(0)=100$, $Z(0)=0$ cell equivalents. Differentiation gives $\dot C=-qC$ and $\dot Z=qC$, while adding the solutions gives $M(t)=100$ at every time. At day 5, about 36.8 equivalents remain viable and 63.2 are necrotic, with no loss of total material. This tests the transfer subproblem at a fixed effective rate; it is not an asserted solution of the whole dynamically coupled oxygen system.

\coursenotebook{14d_coupled_tumour_growth.ipynb}{Notebook 14d} supplies the coupled solver. For a numerical audit, compare the declared discrete accounting and fluxes; a finite-step splitting method need not satisfy an instantaneous continuum identity to machine precision at every substep.
\answerconnection{The rim--core story is a mechanistic hypothesis. The balance specifies which changes could produce a viable deficit and prevents a biologically misleading one-compartment interpretation.}

\solutionslide{20}{Compare states with matched displays}
\biomedicalanswer
\task{Predict the zero-uptake pattern, use trajectories to assess it, and explain why a difference image alone cannot identify a death rate.}
\answerstep{1. Predict the model's limiting behaviour.}
When oxygen uptake is disabled in this normalized model, positive delivery with no-flux oxygen boundaries supports oxygen at its assigned normal baseline throughout active tissue. With the usual zero initial necrotic field and thresholds below baseline, oxygen does not activate death transfer. Proliferation retains its full oxygen modifier and viable growth should match the otherwise identical Fisher--KPP reference.

The prediction belongs to this mathematical closure. It does not say that suppressing consumption in real tissue would eliminate every cause of cell death. Nor is a reduction in the model's delivery factor alone equivalent to oxygen depletion when its consumption term is zero.
\answerstep{2. Describe the check with trajectories.}
In \coursenotebook{14d_coupled_tumour_growth.ipynb}{Notebook 14d}, rerun the declared uptake-zero comparison and inspect minimum active-tissue oxygen, total viable cells and accumulated necrotic material over saved days. Expect baseline oxygen, zero necrotic accumulation from zero initial necrosis and matched viable trajectories within the declared numerical accuracy. Compare equal physical locations and fixed display limits as well as full-domain totals. Differences outside active tissue should not be confused with tissue oxygen.

If the expectation fails, inspect initial necrosis, thresholds, boundary conditions, other parameters and numerical settings before assigning a biological explanation. These instructions define a check rather than report a new run.
\answerstep{3. Explain why a death rate is not identified.}
A smaller viable field can result from slower proliferation, altered transport, death or combinations of these. Remaining necrotic material additionally depends on clearance and previous history. One difference image records a state comparison, not the separate rates that generated it. Trajectories and independent information about the compartments can constrain those rates, but even many observations may leave parameter trade-offs.
\answercheck{The answer should specify initial necrosis, what ``normal oxygen'' means in the lesson and which trajectories test the limit. It should not infer a unique death rate from reduced viable intensity or visible extent.}

\computationalanswer{20}
\task{Explain the assumptions under which zero uptake recovers Fisher--KPP, derive the oxygen result and distinguish the continuum proof from a numerical check.}
\answerstep{1. Substitute zero uptake in the oxygen balance.}
The quasi-steady equation is
\[
0=\ell_o^2\Delta o+p(1-o)-\beta(c/K)o,\qquad
p=[1+\alpha(c+z)/K]^{-1}.
\]
Here $\ell_o$ is a length, $o,p,\alpha,\beta$ are dimensionless and $\Delta$ is the spatial Laplacian. Set $\beta=0$ and $w=o-1$. Since the Laplacian of a constant is zero, the equation becomes
\[
0=\ell_o^2\Delta w-pw.
\]
At a fixed tumour state, $p$ is a prescribed positive spatial function provided the burden and nonnegative suppression parameter are finite and $K>0$.
\answerstep{2. Multiply, integrate and remove the boundary term.}
For constant $\ell_o>0$, multiply by $w$ and integrate:
\[
0=\ell_o^2\int_\Omega w\Delta w\,dV-\int_\Omega pw^2\,dV.
\]
Integration by parts states
\[
\int_\Omega w\Delta w\,dV
=\int_{\partial\Omega}w\,\partial_{\mathbf n}w\,dS
-\int_\Omega|\nabla w|^2\,dV.
\]
No flux for oxygen gives $\partial_{\mathbf n}w=0$, so
\[
0=-\ell_o^2\int_\Omega|\nabla w|^2\,dV-\int_\Omega pw^2\,dV.
\]
The two integrands are nonnegative before their minus signs. A sum of their negatives can be zero only if the integrals vanish. Since $p>0$, the second forces $w=0$ almost everywhere; for regular solutions, $o=1$ throughout tissue. Disconnected tissue components obey the same conclusion under the same positive-delivery assumption.
\answerstep{3. Remove death and recover the logistic term.}
With $0<\theta_d<\theta_p<1$,
\[
g(1)=\min(1,1/\theta_p)=1,\qquad
h(1)=\max(1-1/\theta_d,0)=0.
\]
Therefore $T=0$. If $z_0=0$, then $\partial_tz=-\lambda_z z$ gives $z=0$ for all times. The viable equation reduces to
\[
\partial_tc=\nabla\!\cdot(D\nabla c)+\rho c[1-c/K]_+.
\]
For constant $K>0$, nonnegative $D,\rho$, initial $0\leq c_0\leq K$ and compatible no-flux boundaries, zero and $K$ are lower and upper bounds under the usual comparison assumptions. Hence $0\leq c\leq K$ and the positive-part clipping changes nothing: this is Fisher--KPP.
\answercontinued{20}
\answerstep{4. Explain why the assumptions matter.}
Positive delivery selects the baseline solution; if delivery and uptake both vanish on an isolated component, the no-flux oxygen equation need not select that particular constant. Zero initial necrosis removes the residual crowding contribution; otherwise existing $z$ may still suppress births while clearing. Constant $K$ supports the stated simple upper bound; spatially varying capacity needs a separate argument. The $\ell_o>0$ proof is sufficient, not necessary: if $\ell_o=0$ and $p>0$, the algebraic equation $p(1-o)=0$ also gives $o=1$.
\answerstep{5. Define the numerical comparison.}
Match initial fields, domain, $D$, $\rho$, $K$, grid, saved times and declared solver accuracy. On active voxels, useful dimensionless errors are
\[
e_o=\max_i|o_i-1|,\quad e_z=\max_i|z_i|/K,\quad
e_c=\max_i|c_{\beta=0,i}-c_{r,i}|/K.
\]
Evaluate over the compared times as well as locations. Finite tolerances and discretization must be distinguished from the exact continuum identity. The analytical limit provides an independent target; an overlay alone is not that check.
\answerconnection{The biological prediction becomes a conditional mathematical result. Agreement tests a specified implementation limit; it does not establish that the oxygen closure is physiologically complete.}

\solutionslide{21}{Agents produce different kinds of evidence}
\biomedicalanswer
\task{Follow a source claim to an agent's proposed experiment, identify a contradictory observation and explain what failure would mean biologically.}
\answerstep{1. Preserve the distinction between source and proposal.}
Start with the syllabus's discussion of go-or-grow: a hypothesis distinguishing migratory and proliferative states \citep{hatzikirou2012}. An agent might propose a more specific question: ``Under the chosen oxygen conditions, do viable cells change their movement and division behaviour?'' This is a proposed operationalization for a teaching example, not a claim that the cited study performed the experiment described below.

Trace what the source actually supports, what the agent infers and which additional assumption turns that inference into a testable proposal. The agent's well-written plan is a research artefact, not a new biological observation.
\answerstep{2. Specify observations that address the mechanism.}
A possible laboratory study would compare declared oxygen conditions while following cells over time, recording movement, divisions and survival separately. Track which cells remain observable. If cells disappear because they die or leave the field of view, treating disappearance as increased movement would confuse mechanisms. Conditions, biological preparation, exposure history and analysis rules need to be defined before results are inspected.

Separate biological replicates from repeated analysis of the same recording. A rerun of tracking software can test analysis sensitivity; it does not create another preparation.
\answerstep{3. State what could contradict the proposal.}
If the specified viable population shows no predicted movement change under a comparison capable of detecting the proposed effect, that result challenges this specific prediction in that setting. If the apparent effect instead accompanies substantial death or tracking failure, it challenges the agent's interpretation of movement. Appropriate controls and measurement precision are needed to distinguish absence of evidence from an informative negative result.

Failure would not refute every oxygen-related mechanism or every form of go-or-grow. It would narrow the proposed relationship, motivating changes to the hypothesis, model or experiment. The progression from extraction to hypotheses and empirical knowledge requires new evidence, as emphasized by the agent framework in the syllabus \citep{primeview2026}.
\answercheck{A strong answer gives a falsifiable prediction, measures competing mechanisms and states exactly which claim a failed test would weaken.}

\computationalanswer{21}
\task{Enumerate the four score pairs, calculate selection bias and define an independent evaluation. Explain why reruns are not new biological samples.}
\answerstep{1. Enumerate the sample space.}
Let $Z_1,Z_2$ be independent scores taking $-1$ and $+1$ with equal probability. Each has expectation zero. Independence makes the four pairs equally probable:
\begin{center}
\begin{tabular}{ccc}
\toprule
$Z_1$ & $Z_2$ & Selected maximum\\\midrule
$-1$&$-1$&$-1$\\
$-1$&$+1$&$+1$\\
$+1$&$-1$&$+1$\\
$+1$&$+1$&$+1$\\\bottomrule
\end{tabular}
\end{center}
Therefore $P(\max=-1)=1/4$, $P(\max=+1)=3/4$, and
\[
\mathbb E[\max(Z_1,Z_2)]=(-1)\tfrac14+(+1)\tfrac34=\tfrac12.
\]
Neither candidate has improved; selecting by the observed noisy score changes the expected reported result.
\answerstep{2. Generalize to see the role of search.}
For $N$ independent scores with the same distribution, the maximum is negative only if all $N$ are negative, an event of probability $2^{-N}$. Thus
\[
\mathbb E[\max_{1\leq i\leq N}Z_i]
=-2^{-N}+(1-2^{-N})=1-2^{1-N}.
\]
For $N=1$ the expression is zero; for $N=2$ it is $1/2$; as search expands it approaches one. This is an exact result for the toy noise model, not a universal formula for an agent's performance.
\answerstep{3. Evaluate without reusing the selection noise.}
Let the chosen candidate index be $J$, determined using the first scores. Give each candidate an independent new evaluation score $W_i$, with mean zero and independent of all selection scores. Then
\[
\mathbb E[W_J]=\sum_j P(J=j)\mathbb E[W_j\mid J=j]
=\sum_jP(J=j)\,0=0.
\]
The new evaluation removes the particular noise advantage used for selection under these assumptions. In a real study, reserve evaluation cases, biological samples or outcomes unavailable during candidate development. Merely hiding filenames while reusing the same underlying patient information does not ensure independence.
\answerstep{4. Record the procedure being reproduced.}
An agent trace should include sources, prompts or actions, tool results, candidate-selection rule and stopping criterion. Independent runs on the same dataset can assess variation of that procedure. They cannot change how many independent biological units were observed. A repeated computational analysis and a new biological replication occupy different levels of evidence.
\answerconnection{The biomedical route identifies a test that can contradict the claim. The formal route explains why selecting a convincing run cannot replace that independent test.}

\solutionslide{22}{Place assistance inside the research design}
\biomedicalanswer
\task{Propose a measurement that helps distinguish dispersal from proliferation, and explain why a progression forecast is not a treatment-effect estimate.}
\answerstep{1. Explain the ambiguity in serial appearance.}
A visible region can enlarge because cells move into neighbouring tissue, because local cell numbers increase, or because a thresholded measurement changes. A similar boundary speed may be compatible with different combinations of movement and proliferation. A close fit to a few earlier images therefore need not identify the individual rates.

This ambiguity is especially important when the scan is only indirectly related to cell density. More precise measurement of the same ambiguous boundary may improve reproducibility without completely resolving the biological question.
\answerstep{2. Add information with a different mechanistic connection.}
One possible additional observation is an independently defined tissue measurement of the fraction of cells undergoing division, sampled and interpreted for the question of interest. It may constrain proliferation while spatial observations constrain spread. Another experimental possibility is direct cell tracking over time, with divisions recorded separately. Either requires an explicit account of how the measurement relates to the model's population-level rate.

For imaging, compare spatial density-related profiles or changes at several locations and times rather than only a single volume. Such observations could help distinguish a broader low-density front from local accumulation, but the signal-to-cell relation must be established. A proliferation-associated marker is not automatically the model parameter itself.
\answerstep{3. Separate forecasting from intervention.}
A forecast asks what a later observation is likely to be under the setting represented in the data. A treatment effect asks how outcomes would differ under alternative interventions. Comparing a forecast with what happened supplies only the observed trajectory, not the unobserved trajectory for that same person under another treatment. A causal design or additional justified assumptions are required for the latter question.
\answercheck{The additional measurement should constrain a different aspect of the mechanism, and its observation model should be stated. Good prediction under observed care does not by itself answer which treatment would change that outcome.}

\computationalanswer{22}
\task{Compare persistence and fitted-model residuals at day 90. Explain why a precise grid minimum may leave parameters difficult to distinguish.}
\answerstep{1. Define the forward and observation maps.}
Let $\theta=(s,\rho)$, where $s$ is a dimensionless multiplier of tissue diffusivity and $\rho$ is a rate in day$^{-1}$. The simulated density fraction is $u_\theta=c_\theta/K$. Operator $H$ averages each retained $2\times2\times2$ block:
\[
[Hu]_i=\frac18\sum_{r\in\text{block }i}u_r.
\]
For observation time $t_j$, $y_{j,i}$ is the generated noisy fraction in retained block $i$. Fit days 20, 40 and 60 by minimizing mean squared residual
\[
J(\theta)=\frac1{N_{\rm obs}}\sum_{j,i}
\bigl([Hu_\theta(t_j)]_i-y_{j,i}\bigr)^2
\]
over the declared grid $\Theta_{\rm grid}$. The denominator counts included time--block pairs, not only unique spatial blocks. Day 90 must remain unavailable to the fitting and selection rule.
\answerstep{2. Compare the held-out errors using the same observations.}
For day-90 block $i$, define
\[
r_i^{\rm pers}=y_{60,i}-y_{90,i},\qquad
r_i^{\rm model}=[Hu_{\hat\theta}(90)]_i-y_{90,i}.
\]
Persistence uses the saved noisy day-60 observation, not an inaccessible clean state. Both methods are scored on the same day-90 blocks. For example, if $y_{60}=0.20$, $y_{90}=0.30$ and model prediction is $0.28$, the residuals are $-0.10$ and $-0.02$, with squared errors $0.0100$ and $0.0004$. This single-block illustration does not determine a whole-volume ranking.

Across $n$ retained day-90 blocks, $\mathrm{RMSE}=\sqrt{n^{-1}\sum_i r_i^2}$ is in density-fraction units. Lower is better under this criterion.
\answerstep{3. Read the saved reference result at its proper scope.}
The course's \href{https://github.com/arvidl/brains-to-models/blob/main/paper/figures/growth_extension_values.json}{saved generated-data reference} records generating and recovered parameters $s=1.0$ and $\rho=0.04$ day$^{-1}$. Its day-90 RMSE against noisy observations is approximately $0.004651$ for the model and $0.062877$ for persistence. These are saved results, not a fresh run reported here. The separately stored clean-field RMSE measures a different comparison and should not be substituted for the noisy-observation score.
\answerstep{4. Explain the identifiability limitation.}
A grid search returns the smallest evaluated objective, possibly a unique pair, even if neighbouring parameter pairs produce nearly the same early observations. A flat valley of $J$ means small objective differences can select quite different parameters as noise or model assumptions change. Grid spacing can make the reported pair look more exact than the evidence warrants.

In \coursenotebook{14c_growth_parameter_recovery.ipynb}{Notebook 14c}, the generating pair is included in the search, seed and tissue ratio are known, and generation and fitting share the spatial model. Successful recovery tests this controlled problem. It is neither a patient parameter estimate nor an uncertainty interval. Independent observation types, additional times and sensitivity analysis can help assess which parameter differences are actually distinguishable.
\answerconnection{The biomedical route proposes information that constrains different mechanisms. The objective and held-out comparison make explicit which aspects of recovery and forecasting the current experiment tests.}

\solutionslide{23}{A creative workflow and its provenance}
\biomedicalanswer
\task{Explain oxygen-limited growth in ordinary language and as a visual sequence. Identify a misleading omission and ask a viewer what they understood.}
\answerstep{1. Give a plain-language account.}
``Growing cells use oxygen supplied by their surroundings. If supply does not keep pace, cells may multiply more slowly. Under sufficiently poor conditions in this model, some viable cells become necrotic material. That material can remain for a time and later be cleared. The pattern depends on the supply, consumption and movement assumptions.''

This description distinguishes slower multiplication from death and accumulation from clearance. It does not require equations, and it avoids claiming that every part of a tumour follows the same sequence.
\answerstep{2. Translate the account into a sequence.}
Use matched panels at a few saved days. First show viable density and an oxygen legend. Next show oxygen falling where burden is high, while explaining that some cells remain viable but proliferate more slowly. Then show necrotic material appearing only where the model's death condition is met. Finally show how viable, necrotic and total material can have different trajectories as clearance acts.

Keep the same physical slices, times, field labels and colour limits across comparisons. A changing slice or colour scale should be announced, because either can create an apparent biological change.
\answerstep{3. Identify an omission and test understanding.}
If the sequence jumps from ``less oxygen'' directly to ``all missing viable cells are dead'', it omits reduced proliferation. If necrotic material disappears without explaining clearance, viewers may infer recovery of viable cells. If the model field is drawn over MRI without labelling its origin, viewers may mistake simulation for measured tumour biology.

Ask a viewer: ``Can the viable population be smaller without an equal amount of dead material? What caused the change shown here?'' Their explanation reveals whether the visual story retained the intended distinctions. Emotional impact is another outcome to assess separately. The creative workflow discussed in the syllabus combines human and generated contributions \citep{knudsen2026changed}; its persuasiveness does not replace these comprehension checks.
\answercheck{A strong visual explanation preserves the distinctions among fields and processes, not merely an attractive sequence of colours.}

\computationalanswer{23}
\task{Formulate and test a colour-consistency invariant on saved frames. Distinguish reproducible rendering from replaying an external generation service.}
\answerstep{1. Define a fixed data-to-colour map.}
Let scalar physical value $u$ have declared display limits $l<h$. Normalize and clip by
\[
\eta(u)=\min\!\left(1,\max\!\left(0,\frac{u-l}{h-l}\right)\right),
\qquad \mathrm{colour}(u)=C(\eta(u)),
\]
where $C:[0,1]\to[0,1]^3$ is a fixed red--green--blue map. Clipping defines how values outside the range are handled. The invariant is: the same physical value in any compared frame receives the same mapped RGB triple when $l,h,C$ are unchanged.
\answerstep{2. Verify why frame-specific scaling breaks it.}
For value $u=0.5$, frame limits $[0,1]$ give normalized value $0.5$, while limits $[0,0.5]$ give one. A nonconstant colour map can assign different colours to those two normalized values even though the physical value is unchanged. If $h=l$, the normalization is undefined; a constant-frame convention must be declared instead of dividing by zero.
\answerstep{3. Test two saved frames.}
Record the scalar limits, units, colour map and clipping rule. Select equal-valued probes in both frames, or pass a shared set of known probe values through the saved rendering procedure. Confirm identical mapped triples within the numerical or image-encoding tolerance. Also check a few known unequal values so that an accidentally constant map does not pass as useful scientific rendering.

The invariant concerns the data-to-colour map. A final screen pixel blended with a different background or opacity can look different despite the same mapped RGB value. Match or separately account for those settings, slice coordinates and masks. Equal array indices alone do not establish equal physical locations across different geometries.
\answerstep{4. Preserve the stages of the pipeline.}
Write animation $a=R(d,\phi)$ for saved data $d$ and rendering settings/assets $\phi$. Preserve source arrays, physical geometry, software and asset versions, parameters, transformations and revision decisions. This permits a declared rendering to be reproduced or compared even if a previous external music or image generator no longer returns exactly the same result.

Re-rendering saved assets and regenerating those assets are different procedures. External generation can depend on model version, service state and hidden settings; a seed alone may not specify all of them. Preserve the accepted source assets as part of the scientific communication record.
\answerconnection{The biological story states what the colours are meant to explain. The rendering invariant tests whether the displayed comparison preserves that meaning across frames.}

\solutionslide{24}{A human response to AI change}
\biomedicalanswer
\task{Describe an activity to delegate and one where participation has meaning. Name an observation that could change the view.}
\answerstep{1. Make the choice specific to a purpose.}
A defensible sample choice is to delegate an initial formatting task: organizing already verified measurements and dates into a draft table. Its output is inspectable against sources. Retain personal participation in discussing an uncertain finding with someone whose concerns, priorities and understanding need to be heard. That interaction involves listening and adapting the explanation, rather than simply delivering fluent text.

This is a sample value judgment, not a rule that every person must make the same division. The film discussed in the syllabus expresses human ambivalence about AI-assisted change \citep{knudsen2026changed}. Enthusiasm, fear and perceived efficiency can coexist without determining which workflow serves a particular person.
\answerstep{2. Explain the trade-off concretely.}
Delegating formatting could free attention, or it could create checking work if the draft mismatches dates. Participating in a conversation may matter because a recipient needs room to ask what uncertainty means for them. A professional's excitement about a faster tool does not automatically answer that need. The relevant comparison includes the accepted output and the experience of people involved.
\answerstep{3. State observations that could revise the choice.}
If the delegated table repeatedly introduces consequential errors or takes longer to verify than to prepare directly, revise the delegation. If a particular communication aid improves the recipient's demonstrated understanding and gives them more opportunity to ask questions, incorporate it while reassessing the role of personal participation. Conversely, if recipients prefer and understand another arrangement, that evidence should challenge the original assumption about what helps them.

Measure work time, error correction, comprehension and reported experience separately. A strong emotional response to the film or tool can be meaningful without demonstrating factual learning.
\answercheck{The answer should state whose purpose is served, what is being delegated and a concrete observation that could change the choice. Personal meaning and measurable task performance are related questions, not interchangeable outcomes.}

\computationalanswer{24}
\task{Define comprehension and emotional-response scores for assigned film and transcript groups. Compare them separately and interpret the total-time example.}
\answerstep{1. Define two non-interchangeable outcomes.}
Let $C_i$ count correct answers to five fixed comprehension questions for participant $i$, so $C_i\in\{0,1,\ldots,5\}$. Define each correct answer before observing group responses, including the key uncertainty distinctions. Let $E_i\in\{0,1,\ldots,10\}$ be a declared self-report of emotional arousal from ``none'' to ``very strong''. This measures response intensity, not whether the feeling is pleasant, and does not directly measure knowledge.

Keep film and transcript factual content comparable for the declared question. If sound and imagery are part of the intervention, the comparison estimates their combined presentation effect rather than identifying one component automatically.
\answerstep{2. Calculate separate group contrasts.}
For an invented example with four participants per group, suppose
\[
\begin{array}{c|cc}
& C & E\\\hline
\text{Film}&(3,4,2,3)&(8,7,9,8)\\
\text{Transcript}&(4,3,3,2)&(3,4,2,3)
\end{array}
\]
The mean comprehension scores are both $3$, whereas mean arousal is $8$ for film and $3$ for transcript. Thus
\[
\widehat\Delta_C=3-3=0,\qquad
\widehat\Delta_E=8-3=5.
\]
The outcomes even have different scales. The emotional contrast cannot be read as a five-point comprehension improvement. The small invented sample also supplies no claim of statistical precision or population equivalence of comprehension.
\answerstep{3. Explain the design and inference.}
Random assignment, a consistent intervention, comparable assessment timing and an appropriate analysis support a group contrast. Prior knowledge, missing responses and scoring should be handled explicitly. Blind comprehension scoring where possible; use new questions or another suitable design to avoid learning from repeated testing. A difference in $E$ alone provides no logical implication about $C$; both must be measured.
\answerstep{4. Calculate the complete workflow time.}
For generation, review and revision times in minutes, $T=T_g+T_r+T_v$. The syllabus's invented comparison is
\[
T_{\rm initial}=5+1+1=7,\qquad
T_{\rm assisted}=1+5+3=9.
\]
Generation is four minutes faster, but total time is two minutes greater, about $28.6\%$ of the original total. This calculation assumes non-overlapping, comparably defined tasks and the same acceptance standard. Distinguish person-work time from machine waiting and elapsed time; overlapping activities cannot simply be added to obtain elapsed time.
\answerconnection{Human purposes identify which effects matter. Separate outcome and time definitions make it possible to assess those effects without substituting excitement or first-draft speed for understanding and accepted work.}

\solutionslide{25}{Synthesis and a next experiment}
\biomedicalanswer
\task{Explain a positive result without beginning with a formula. Use a teach-back question and name an observation that could change the interpretation.}
\answerstep{1. Count people in the stipulated population.}
Imagine 10,000 people, of whom 100 have the condition and 9,900 do not. The test finds 90 of the 100 affected people. It also produces positive results in 990 of the unaffected people. Altogether 1,080 people test positive, and 90 of them have the condition.

One possible explanation is: ``In this example, about one person in twelve with a positive result has the condition. The positive result deserves interpretation, but it does not by itself confirm the condition. The condition is uncommon in this population, so positive results from the much larger unaffected group make up most of the positives.''

These are expected counts based on hypothetical prevalence and test characteristics, not the measured outcome of screening a real group. They describe a specified population; an individual's full information may support another probability.
\answerstep{2. Check understanding in ordinary language.}
Ask: ``Of the people who tested positive in this example, are most affected or unaffected? Why?'' A useful answer recognizes that most positives come from the large unaffected group, while also recognizing that some positives are true. ``The test is 90\% sensitive, so a positive means a 90\% chance of disease'' reverses the question being asked.

Another misunderstanding is ``a positive is usually false, so it can always be ignored''. The next action also depends on consequences, further evidence and the context in which testing was used.
\answerstep{3. Name information that changes the account.}
A different condition frequency in the tested population changes the composition of positive results. Different test performance or additional well-characterized evidence can also change interpretation. Explain what is measured and how it relates to the first result; a second test should not be assumed independent merely because it has another name.
\answercheck{Preserve the denominator: the question concerns people with positive results. The explanation should keep uncertainty and action separate while identifying what new observation would matter.}

\computationalanswer{25}
\task{Derive PPV and the cost-based action boundary. Distinguish prevalence, sensitivity and the event probability used for a decision.}
\answerstep{1. Construct the complete expected-count table.}
Let $D$ denote disease and $+$ a positive test. With $\pi=P(D)=0.01$ and $\mathrm{Se}=P(+\mid D)=0.9$, $\mathrm{Sp}=P(-\mid\neg D)=0.9$, the expected counts are
\begin{center}
\begin{tabular}{lrrr}
\toprule
& Disease & No disease & Total\\\midrule
Positive & 90 & 990 & 1,080\\
Negative & 10 & 8,910 & 8,920\\\midrule
Total & 100 & 9,900 & 10,000\\\bottomrule
\end{tabular}
\end{center}
Checking row and column totals catches many denominator errors.
\answerstep{2. Derive rather than memorize the fraction.}
Conditional probability gives $P(D\mid+)=P(D\cap+)/P(+)$. The joint true-positive probability is
\[
P(D\cap+)=P(+\mid D)P(D)=\mathrm{Se}\,\pi.
\]
Disease and no disease partition the population, so total positive probability is
\[
P(+)=\mathrm{Se}\,\pi+(1-\mathrm{Sp})(1-\pi).
\]
Combining these expressions yields
\[
\mathrm{PPV}=\frac{\mathrm{Se}\,\pi}
{\mathrm{Se}\,\pi+(1-\mathrm{Sp})(1-\pi)}
=\frac{0.009}{0.009+0.099}=\frac1{12}\approx0.0833.
\]
Sensitivity conditions on disease; PPV conditions on a positive result. They have different denominators and need not be equal. At prevalence $0.10$ with the same test characteristics, PPV becomes $0.09/(0.09+0.09)=0.5$.
\answercontinued{25}
\answerstep{3. Derive the decision boundary independently.}
For event probability $p$ given available decision evidence, stipulate false-positive cost 1, false-negative cost 4 and zero loss for correct actions. The two expected losses are $1-p$ and $4p$, so
\[
1-p<4p\quad\Longleftrightarrow\quad p>0.2.
\]
At equality, both actions cost $0.8$ in expectation. If a positive test is the only new evidence in the stipulated population and no other information is used, set $p=\mathrm{PPV}=1/12$. Acting then costs $11/12$ and waiting $4/12=1/3$ in this toy loss model. The preferred toy action follows from the supplied consequences, not from PPV alone. Other costs or actions change the decision problem.
\answerstep{4. Check definitions at the boundaries.}
Prevalence $\pi$ is the population probability before conditioning on this test. Decision probability $p$ uses the available evidence; equating it with PPV requires the declared conditioning information. If $P(+)=0$, PPV is undefined because the conditioning event cannot occur. That differs from PPV zero with a positive denominator, where positives occur but none are true positives in the model.

\coursenotebook{15_ppv_and_decisions.ipynb}{Notebook 15} permits separate exploration of counts, conditional probabilities and stipulated decision costs. Changing population size rescales expected counts while leaving PPV fixed when prevalence and test characteristics stay fixed.
\answerconnection{Population counts communicate what a positive result means. Conditional probabilities establish the denominator, and a separate loss model makes the action assumptions explicit.}

\clearpage
\gdef\currentsolutionheading{Sources and follow-up}
\section*{Sources and reproducible follow-up}
\addcontentsline{toc}{section}{Sources and reproducible follow-up}
The exercise wording and notation follow the separate syllabus. References below support the concepts discussed here; invented examples are not estimates from those publications. The three documents share formatting and one bibliography database.

For notebook exploration, record the notebook and source revision, parameter values, seed, grid, time settings and the quantity measured. Keep a verbal prediction separate from the observed result. Compare the result with the mathematical check and then explain which biological question remains open. Saved reference values cited for slide 22 come from \href{https://github.com/arvidl/brains-to-models/blob/main/paper/figures/growth_extension_values.json}{the course's generated-data results file}; they are not patient measurements or a new execution reported here.

\textbf{Reviewing an answer.} A biomedical answer should connect a mechanism or anatomical relationship to an observable consequence, consider a plausible alternative and name evidence that distinguishes them. A computational answer should define its quantities, justify each operation, check units or a limiting case and state the assumptions required by its conclusion. Agreement between the paths is strongest when both refer to the same observable quantity.

\clearpage\phantomsection
\gdef\currentsolutionheading{References}
\addcontentsline{toc}{section}{References}
% Each document lists only the sources it cites, from one database.
\begingroup
\small
\setlength{\bibsep}{3pt}
\interlinepenalty=10000
\bibliographystyle{plainnat}
\bibliography{references}
\endgroup

\end{document}
